% ======================================================================
% paper/companion/removed-sections-draft.tex  (Phase 06 working copy)
% Verbatim archival copy of the pre-restructure manuscript front half.
% Source: paper/manuscript.tex @ commit 4f1815a (lines 129-549).
% NOT for compilation: unformatted draft; P10 formats the companion.
% ======================================================================
%
% §1 Introduction and §2 Related Work are rewritten wholesale per the
% P05 blueprint (§7); this file preserves every word of the originals,
% including content cut from the narrow paper (faculty-alignment table,
% four-version history, scope qualifiers, continual learning, causal
% representation learning, cognitive evaluation, five-gap map, per-paper
% differentiations).
%
\section{Introduction}\label{sec:introduction}
% ============================================================

The standard framing of systematic compositionality failure is precise
and reproducible. Agents trained to high in-distribution accuracy on
sequence and language tasks fail near-completely when test distributions
require zero-shot recombination of primitives trained in isolation.
\cite{lake2018generalizationwithoutsystematicity} demonstrated this on
SCAN --- seq2seq models achieve above 99\% accuracy on random test splits
and below 2\% on the add-primitive split. % CLAIM:C-001 The broader benchmark
literature is consistent: COGS \citep{kim2020cogscompositionalgeneralizationchallenge}
finds 96--99\% in-distribution accuracy alongside 16--35\% compositional
generalisation accuracy; % CLAIM:C-002 CFQ \citep{keysers2020measuringcompositionalgeneralizationcomprehensive}
shows performance degradation growing monotonically with compound
divergence; PCFG-SET \citep{hupkes2020compositionalitydecomposed} documents
failures on productivity, systematicity, and substitutivity tests across
RNNs, CNNs, and Transformers.

Why does this failure pattern repeat across every benchmark? The answer
lies in what training pipelines actually optimise. They have one formal
target: minimise expected loss over the training distribution. This
objective efficiently grows one variable --- \textbf{parametric depth}
$\delta_A(d,t)$ --- but has no formal mechanism for increasing a
second, independent variable: \textbf{schema coherence} $\sigma_A(d,t)$,
the degree to which representations are restructured around deep
governing principles.

The $\Sigma$-Model has evolved through four versions to address this
asymmetry. V1.0 formalised the $\delta_A$/$\sigma_A$ distinction. V2.0
extended it to Attention, Executive Functions, Metacognition, and Social
Cognition --- each with formal variables, ODEs, and benchmark protocols.
V3.0 extended all variables to a domain $\times$ modality product space
and formalised benchmark validity as a measurable object. V3.0$+$ adds a
reliability function and three practical protocols that close the remaining
measurement gaps.

\subsection{Central Claim}\label{sec:central-claim}

\begin{quote}
Training pipelines that optimise $\delta_A(d,t)$ without formally
targeting $\sigma_A(d,t)$ are predicted to systematically produce
agents that pass in-distribution evaluation while failing
out-of-distribution recombination --- not because they lack depth,
but because their training regimes suppress the schema crystallisation
that converts parametric depth into principled generalisation capacity.
Furthermore, the suppression mechanism extends to attentional allocation
($\alpha_A$), executive control ($\Xi_A$), metacognitive self-modelling
($\hat{M}_A$), and inter-agent schema communication ($\mu_{AB}$) ---
all of which interact multiplicatively with $\sigma_A$ through the
formal coupling terms derived below.
\end{quote}

\paragraph{Scope qualification.}
Our central claim applies to agents trained on compositional
generalisation benchmarks (SCAN, COGS, PCFG-SET) using standard
gradient-based optimisation. We do \emph{not} claim that
$\sigma_A$-targeting is the \textbf{only} path to compositional
success. Large-scale pretraining \citep{kaplan2020scaling}, whose scaling laws concern perplexity rather than compositional generalisation specifically,
in-context learning \citep{min2022rethinking}, whose mechanisms are partly explained by label-space and format conditioning rather than input-label mapping, and dual-process
architectures \citep{noviello2025neuroscienceinspiredualprocess} may achieve similar ends via
different mechanisms. The Σ-Model framework provides one
\emph{sufficient} mechanism (schema-coherence cultivation) and
predicts its empirical signatures; it does not claim \emph{necessary}
exclusivity.

\subsection{Cognitive Faculty Alignment}\label{sec:cognitive-faculty-alignment}

Five cognitive faculties identified as evaluation gaps
by~\cite{burnell2026measuringprogressagi} map directly onto
$\Sigma$-Model state variables. This mapping grounds each faculty in a
measurable dynamical variable:

\begin{table}[htbp]
\centering
\caption{\textbf{Cognitive faculty alignment with Σ-Model variables.}}
\label{tab:faculty-alignment}
\begin{tabular}{>{\raggedright}p{2.8cm} >{\raggedright}p{3.2cm} >{\raggedright\arraybackslash}p{7.2cm}}
\toprule
\textbf{Faculty} & \textbf{Primary Variable} & \textbf{Mechanism} \\
\midrule
\textbf{Learning} & $\sigma_A(d,t)$, $\delta_A(d,t)$ & OOD gap as schema proxy; compositional generalisation \\
\textbf{Metacognition} & $\hat{M}_A(d,t)$, $\zeta_A(d,t)$ & Self-model accuracy; calibration error dynamics \\
\textbf{Attention} & $\alpha_A(d,t)$, $C_A(d,t)$ & Attentional fidelity to generative structure \\
\textbf{Executive Functions} & $\Xi_A^P$, $\Xi_A^I$, $\Xi_A^F$ & Planning, inhibition, cognitive flexibility \\
\textbf{Social Cognition} & $\mu_{AB}$, $\tau_A(B,d,t)$, $\Sigma_{A,B}$ & Schema legibility, theory of mind, collective field \\
\bottomrule
\end{tabular}
\end{table}

\subsection{Argument Roadmap}\label{sec:argument-roadmap}

The central claim (Section~\ref{sec:central-claim}) is established through the
following proof structure, with each step mapped to a paper section:

\textbf{THESIS:} $\delta$-only optimisation $\Rightarrow$ $\sigma$-suppression $\Rightarrow$ OOD failure

\textbf{PROOF STRUCTURE:}
\begin{itemize}[nosep]
\item \textbf{Step 1 (Section~\ref{sec:core-framework}):} Define observables ($\delta_A$, $\beta_A$, $\sigma_A$) and show $\sigma_A$ is latent
  \begin{itemize}[nosep]
  \item Propositions~3.6--3.7: $\sigma_A$ cannot be inferred from $\delta_A$ alone % CLAIM:C-021
  \end{itemize}
\item \textbf{Step 2 (Section~\ref{sec:cognitive-dimensions}):} Show standard SGD dynamics suppress $\sigma_A$
  \begin{itemize}[nosep]
  \item Theorem~4.1: $\partial\sigma_A/\partial t|_{\text{SGD}} < 0$ when $\Omega_{SL}$ is high \eqref{eq:sigma-ode}
  \end{itemize}
\item \textbf{Step 3 (Section~\ref{sec:phase-structure}):} Characterise failure mode phase structure
  \begin{itemize}[nosep]
  \item Proposition~4.1--4.4: Σ-Model predicts Phase~2 entry is necessary and sufficient for OOD success
  \end{itemize}
\item \textbf{Step 4 (Section~\ref{sec:predictions}):} Empirical signatures detectable in benchmarks
  \begin{itemize}[nosep]
  \item Predictions~1--9: Each maps to one aspect of thesis [Section~\ref{sec:statistical-framework}]
  \end{itemize}
\end{itemize}

The benchmark specification and target validity
components are provided in Section~\ref{sec:empirical}.

\paragraph*{Contributions.}
This paper makes three principal contributions:
\begin{enumerate}[nosep]
\item \textbf{A dynamical theory of $\sigma$-suppression.}
  Standard SGD creates a stable low-$\sigma_A$ equilibrium --- the
  ``$\sigma$-trap'' --- explaining why high-ID/low-OOD agents arise
  systematically, not as implementation bugs
  (Theorem~4.1, Proposition~4.2).
\item \textbf{A formal proxy architecture for latent $\sigma_A$.}
  A two-tier estimation protocol (\S\ref{sec:online-estimation}) with
  established convergence bounds (Propositions~3.6--3.7) makes $\sigma_A$
  measurable without requiring ground-truth access.
\item \textbf{A falsifiable prediction framework.}
  Eight core predictions and one meta-prediction (\S\ref{sec:predictions}),
  each anchored to a specific phase transition and distinguishable from
  depth-only accounts of generalisation failure.
\end{enumerate}

\paragraph{The delta from prior art.}
Existing frameworks attribute compositional failure to architectural
constraints (MIRAGE's dual-process requirement), capacity limits
(RASP's expressivity bounds), or data distribution issues (mutual
exclusivity, dataset cartography). Each targets a different variable
--- dual-process architecture, attention span, lexical overlap --- but
none identifies the shared dynamical mechanism that makes these
failures systematic rather than contingent. The $\Sigma$-Model
proposes that mechanism: $\sigma_A$-suppression is a bifurcation
phenomenon, not a capacity or data problem. Standard training creates
a stable low-$\sigma_A$ equilibrium; all existing interventions succeed
only insofar as they incidentally push $\sigma_A$ above
$\sigma_{\text{critical}}$. This reframes compositional failure from a
symptom to be patched into a dynamical system to be navigated.

The thesis above depends on a critical premise: that $\sigma_A$ is
absent from every existing framework that touches on compositional
generalisation --- that the field has built extensive machinery for
optimising $\delta_A$ while leaving $\sigma_A$ unmodelled and
unmeasured. We now verify this premise across five research areas.
Each area addresses a $\Sigma$-Model variable tangentially, but none
formalises schema coherence as a dynamical variable with its own ODE,
suppression mechanism, or phase transition threshold. Making this gap
precise is the purpose of the review below.

\noindent\rule{\linewidth}{0.4pt}

% ============================================================
\section{Related Work and Gap Analysis}\label{sec:related-work}
% ============================================================

The Σ-Model occupies a specific gap: none of the five research areas below models the dynamics of schema coherence. We review each to make this gap precise.

\subsection{Curriculum Learning}\label{sec:curriculum-learning}

\cite{bengio2009curriculumlearning} established the foundational result:
easy-to-hard sample ordering improves generalisation by providing
representational scaffolding. Subsequent work elaborated difficulty
measurement \citep{kumar2010selfpacedlearning,portelas2020automaticcurriculumlearningdeep,
narvekar2020curriculumlearningreinforcementlearning} and pacing schedules.
Every variant maximises the rate at which $\delta_A$ increases.
Recent empirical work~\cite{mordig2026rethinkingeasytohard} confirms
that easy-to-hard ordering provides no robust advantage for
deductive reasoning benchmarks, reinforcing the gap between
depth-growth optimisation and schema crystallisation.

\textbf{Gap statement.} Curriculum learning provides training-sequence
optimisation for depth growth rate but has no formal account of
$\sigma_A(d,t)$ dynamics, no triggering conditions for schema
crystallisation, and no prescription for the qualitative shift in
training regime that $\sigma_{\text{critical}}$-crossing requires. Σ-Model
is not a curriculum method: it formalises the \emph{phase structure} that
curriculum methods navigate without modelling. The
$\sigma_A$ ODE~\eqref{eq:sigma-ode} specifies what
$\sigma_{\text{critical}}$-crossing requires; curriculum learning
addresses only $\delta_A$-growth rate within a single phase.

\subsection{Compositional Generalisation}\label{sec:compositional-generalisation}

The SCAN/COGS/CFQ/PCFG-SET benchmark family collectively documents the
high-$\delta_A$/low-$\sigma_A$ failure mode empirically~\citep{mondorf2026compositionalarc}.
The consistent diagnosis: models encode statistical regularities rather than the
compositional rules governing the distribution.

\textbf{Gap statement.} The compositional generalisation literature
characterises the failure mode and provides measurement proxies for
$\sigma_{\text{critical}}$-crossing, but does not formalise $\sigma_A$
as the training variable whose dynamics the failure reveals, does not
identify what suppresses it, and does not specify how to design training
regimes that reliably induce schema crystallisation.

\paragraph{Recent Work Differentiation}

\textbf{\cite{noviello2025neuroscienceinspiredualprocess} --- MIRAGE dual-process model.}
Noviello et al.\ propose a dual-process model in which compositional
generalisation is supported by a separate symbolic System-2 module. The Σ-Model framework demonstrates that
this structural postulate is unnecessary by formalising compositional
capacity as a developmental trajectory within a single ODE system.
Specifically, schema coherence $\sigma_A(d,t)$ grows continuously
through the $\sigma_A$ ODE \eqref{eq:sigma-ode}, where the
$\alpha_A$ gating term directs neural training effort toward
compositional structure without requiring a separate symbolic substrate.
The $\sigma_{\text{critical}}$ bifurcation \eqref{eq:sigma-critical}
defines a threshold crossing that is achievable entirely within the
neural parameter space --- no symbolic overlay is invoked.
Thus Σ-Model predicts that MIRAGE's success is bounded by its incidental
$\sigma_A$ elevation, not by the dual-process architecture itself.

\textbf{\cite{li2023slogstructuralgeneralizationbenchmark} --- SLOG benchmark recursion failures.}
Li et al.\ document persistent failures on unseen recursion depths.
Σ-Model predicts these failures from the $\sigma_A$ ODE as a Phase~1
diagnostic: agents with $\sigma_A \approx 0$ cannot transfer recursive
patterns to unseen depths because the schema growth term is gated by
$\alpha_A$, which is suppressed when surface-reward pressure
$R_A^{\text{surface}}$ is high.
Thus Σ-Model predicts that recursion failures are a
$\sigma_A$-diagnostic, not a capacity limit.

\textbf{\cite{bruns2025exploringcompositionalgeneralizationin} --- RASP structural representability.}
Bruns demonstrates via RASP that transformers possess the structural capacity
to represent compositional operations. Σ-Model's framework is
complementary: the RASP proof establishes that
$\sigma_{\text{critical}}$-crossing is theoretically possible, while
Σ-Model's ODE system explains why standard training does not achieve it.
Thus Σ-Model distinguishes capacity from utilisation --- a distinction
RASP alone cannot address.

\textbf{\cite{jiang2022mutualexclusivity} --- Mutual exclusivity training on COGS.}
Jiang et al.\ demonstrate that mutual exclusivity (ME) training pushes
COGS lexical scores substantially higher while structural
compositionality splits remain below 1\%. % CLAIM:C-003 Σ-Model's $\sigma_A$ ODE
predicts this dissociation precisely: ME training increases
$\delta_A$ at the lexical level but does not gate through $\alpha_A$
for structural compositional regularities, so
$\rho \cdot P_A \cdot \alpha_A \cdot (1 - \sigma_A) \approx 0$
regardless of $P_A$ magnitude.
Thus Σ-Model predicts that lexical remediation leaves structural
compositionality gaps intact --- an empirically verifiable dissociation.

\textbf{\cite{swayamdipta2020datasetcartography} --- Dataset cartography curricula.}
Swayamdipta et al.\ show that dataset cartography yields performance
gains without closing structural compositionality gaps. This result is
predicted by Σ-Model's gap statement: curriculum learning optimises
$\eta(d,t)$ but has no formal mechanism for increasing $\sigma_A$.
Thus Σ-Model predicts that map-based curricula will improve ID
accuracy without closing the OOD gap.

\textbf{\cite{redhardt2025} --- Scaling compositional generalisation (NeurIPS Spotlight).}
Redhardt et al.~show that scaling can lead to compositional generalisation,
potentially challenging Σ-Model's central claim that structural gaps persist at
scale. Within Σ-Model, however, scaling-driven gains are predicted to reflect a
Phase~1$\to$2 transition driven by increased $\delta_A$ that eventually crosses
$\sigma_{\text{critical}}$, but at suboptimal resource efficiency compared with
targeted $\alpha_A$ interventions \eqref{eq:alpha-ode}. Scaling alone does not
address the $\alpha_A$ gating mechanism, predicting slower $\sigma_A$ growth per
unit of training compute.
Thus Σ-Model predicts that scaling will eventually close the OOD gap, but at
substantially larger compute budgets than targeted $\sigma_A$ interventions.

\textbf{\cite{mondorf2026compositionalarc} --- Compositional-ARC meta-learning.}
Mondorf et al.~approach compositionality in spatial reasoning through meta-learning
over procedurally generated ARC tasks. The $\Sigma$-Model makes a sharp prediction:
meta-learning enables $\sigma_A$ to crystallise \emph{per task} by compressing the
learning trajectory into few-shot adaptation, but the resulting schemas are
task-specific rather than domain-general. Thus Compositional-ARC should succeed on
ARC-homologous compositions while transferring poorly to structurally different
domains --- a dissociation that $\delta_A$-only accounts cannot explain.

\subsection{Continual Learning and Decay}\label{sec:continual-learning}

\cite{mccloskey1989catastrophicinterference},
\cite{goodfellow2013empiricalinvestigationcatastrophicforgetting}, and
\cite{kirkpatrick2017overcomingcatastrophicforgetting} address parametric
overwriting under sequential training. Every mitigation strategy targets
preservation or recovery of parametric content.

\textbf{Gap statement.} The continual learning literature formally
accounts for parametric decay $\lambda_c$ under sequential training but
has no formal object for the domain frontier $\Delta(d,t)$ and no
mechanism for representing frontier obsolescence $\lambda_f(d,t)$ as a
distinct decay process with different intervention implications.

\textbf{Connection to schema coherence.} Parametric decay $\lambda_c$ in
continual learning is the temporal analogue of schema coherence suppression
$\sigma_A$ in compositional generalisation. Just as $\lambda_c$ causes
forgetting of previous tasks, $\sigma_A$-suppression causes ``forgetting''
of compositional structure. The Σ-Model framework connects these through the
schema-mediated decay reduction mechanism \eqref{eq:effective-decay}:
schema coherence acts as a protective factor against both forms of decay,
with $\gamma_\sigma$ modulating the degree of protection.

\subsection{Causal and Structured Representation Learning}\label{sec:causal-repr}

\cite{scholkopf2021towardcausalrepresentationlearning} establish that
causal representations support OOD generalisation.
\cite{mohan2024structureindrlsurveyopenproblems} document seven structural
prior incorporation patterns. \cite{torresan2024disentangledrepresentationscausalcognition}
distinguish weak from strong disentanglement.

\textbf{Gap statement.} Causal and structured representation learning
specify the target state corresponding to high $\sigma_A(d,t)$ but do
not formalise $\sigma_A$ as a dynamic scalar with its own ODE, do not
model the training process that builds it or the shortcut-learning pressure
$\Omega_{SL}(d,t)$ \citep{geirhos2020shortcutlearning} that suppresses it, and do not connect it to
attentional allocation, executive control, metacognitive self-modelling,
or collective schema communication.

\subsection{Cognitive Evaluation of AI Systems}\label{sec:cognitive-eval}

\cite{burnell2026measuringprogressagi} propose a Cognitive Taxonomy of
ten faculties and identify Learning, Metacognition, Attention, Executive
Functions, and Social Cognition as having large evaluation coverage gaps.
\cite{chollet2019measureintelligence} proposes ARC as a generalisation
benchmark. \cite{morris2024levelsagi} establish a Levels of AGI framework.

\textbf{Gap statement.} Existing cognitive evaluation frameworks specify
what should be measured but provide no formal theoretical grounding for
why specific task designs isolate specific faculties. The Σ-Model Benchmark
Protocol (Section~\ref{sec:benchmark-protocol}) provides that grounding
by deriving benchmark designs directly from formal variable structures.

\subsection{Synthesis --- The Five-Gap Map}\label{sec:five-gap-map}

The five gaps reduce to a single pattern: each literature addresses
one variable while missing $\sigma_A$. The following table consolidates
this diagnosis.

\begin{center}
\renewcommand{\arraystretch}{1.3}
\begin{tabular}{>{\raggedright}p{3.5cm} >{\raggedright}p{3.5cm} >{\raggedright\arraybackslash}p{6.5cm}}
\toprule
\textbf{Literature} & \textbf{Σ-Model Variable Addressed} & \textbf{Σ-Model Variable Missing} \\
\midrule
Curriculum Learning & $\delta_A$ growth rate & $\sigma_A$ dynamics, $\alpha_A$, $\Xi_A$ \\
Compositional Generalisation & $\sigma_A$ failure mode (empirical) & $\sigma_A$ formation mechanism, suppression \\
Continual Learning & $\lambda_c$ (parametric decay) & $\lambda_f$ (frontier obsolescence), $\sigma_A$ coupling \\
Causal/Structured Repr. & $\sigma_A$ target state & $\sigma_A$ developmental trajectory, $\hat{M}_A$, $\mu_{AB}$ \\
Cognitive Evaluation & Faculty identification & Formal theoretical grounding for benchmark design \\
\bottomrule
\end{tabular}
\end{center}

\paragraph{Comparison with Contemporary Frameworks.}
Table~\ref{tab:framework-comparison} positions Σ-Model relative to recent
approaches addressing compositional generalisation failure.

\begin{table}[htbp]
\centering
\caption{\textbf{Comparison of compositional generalisation frameworks.}}
\label{tab:framework-comparison}
\small
\begin{tabular}{>{\raggedright\arraybackslash}p{2.2cm}>{\raggedright\arraybackslash}p{2.5cm}>{\raggedright\arraybackslash}p{2.5cm}>{\raggedright\arraybackslash}p{4.8cm}}
\toprule
\textbf{Framework} & \textbf{Failure Mode} & \textbf{Mechanism} & \textbf{Relation to $\Sigma$-Model} \\
\midrule
MIRAGE \citep{noviello2025neuroscienceinspiredualprocess}
  & System-1/System-2 competition
  & Dual-process architecture
  & Formalises System-2 ($\sigma_A$) growth dynamics \\

RASP \citep{bruns2025exploringcompositionalgeneralizationin}
  & Capacity gap
  & Restricted access sequences
  & Predicts \emph{when} capacity is used ($\sigma_A$-dependent) \\

Mutual Exclusivity \citep{jiang2022mutualexclusivity}
  & Lexical overlap interference
  & Augmentation-induced sparsity
  & ME targets $\sigma_A$ via data; $\Sigma$-Model targets $\sigma_A$ directly\\

Dataset Cartography \citep{swayamdipta2020datasetcartography}
  & Example difficulty mismatch
  & Curriculum ordering
  & Extends to $\sigma_A$-dependent phase curriculum \\

\textbf{$\Sigma$-Model (Ours)}
  & \textbf{$\sigma_A$-suppression trap}
  & \textbf{Coupled ODE + phase curriculum}
  & \textbf{Predicts conditions where prior methods succeed/fail} \\

Compositional-ARC \citep{mondorf2026compositionalarc}
  & Cross-domain meta-generalisation gap
  & Meta-learning + spatial reasoning
  & Meta-learns per-task schemas; $\Sigma$-Model predicts limited cross-domain transfer \\
\bottomrule
\end{tabular}
\end{table}

\paragraph{Distinctive contribution.}
Unlike frameworks that propose \emph{architectural} solutions (MIRAGE,
RASP), \emph{data-centric} solutions (Mutual Exclusivity, Dataset
Cartography), or \emph{meta-learning} solutions (Compositional-ARC),
Σ-Model proposes a \emph{dynamical-systems} explanation:
compositional failure is not merely a capacity or data problem, but a
\textbf{bifurcation phenomenon}. Standard training creates a stable
low-$\sigma_A$ equilibrium (Proposition~4.2); escape requires crossing
a critical threshold via phase-aware curriculum. This predicts that
architectural and data-centric methods will succeed only when they
incidentally raise $\sigma_A$ above $\sigma_{\text{critical}}$---a
testable meta-prediction (Meta-Prediction~M1; Section~\ref{sec:meta-M1}).

\noindent\rule{\linewidth}{0.4pt}

% ============================================================

% ======================================================================
% ======================================================================
% §3 The Σ-Model — Core Framework (original, manuscript lines 316–991 @ post-§1/§2 rewrite)
% Preserved verbatim: breadth β_A, intersection activation Ψ_A, shortcut
% threshold D*, proxy/implementation architecture, extensions, and the
% timescale-separation parts — content moved to the companion per the P05
% blueprint (§7); the narrow paper keeps the two-variable core.
% ======================================================================
\section{The Σ-Model --- Core Framework}\label{sec:core-framework}
% ============================================================

The Σ-Model formalises agent knowledge development as a coupled
dynamical system (Figure~\ref{fig:sigma-framework}). This section presents the foundational architecture
--- the three core state variables (depth, breadth, schema coherence),
their ODEs, intersection activation, and the shortcut threshold.
Subsequent sections extend this core with additional cognitive
dimensions (Section~\ref{sec:cognitive-dimensions}), multimodal coverage
and benchmark validity (Section~\ref{sec:multimodal}), and reliability
protocols (Section~\ref{sec:reliability}), all of which couple back into
the system established here.

\subsection{The Three Core Knowledge Dimensions}\label{sec:three-core}

Let every domain of knowledge be indexed by $d \in D$, where $D$ is the
set of all knowledge domains. For a given agent $A$ at time $t$,
we track three orthogonal quantities per domain.

\subsubsection{Depth \texorpdfstring{$\delta_A(d,t)$}{delta\_A(d,t)}}\label{sec:depth}

The structural complexity and principled coherence of the agent's
internal representation of domain $d$. Not raw parameter volume ---
the degree to which knowledge is organised around deep causal principles
rather than surface-statistical features. Bounded above by the domain
frontier $\Delta(d,t)$:
\begin{equation}\label{eq:depth-bounds}
0 \leq \delta_A(d,t) \leq \Delta(d,t)
\end{equation}

\textbf{Relative depth:}
\begin{equation}\label{eq:relative-depth}
\delta_A^{\text{relative}}(d,t) = \frac{\delta_A(d,t)}{\Delta(d,t)} \in [0,1]
\end{equation}

\subsubsection{Breadth \texorpdfstring{$\beta_A(d,t)$}{beta\_A(d,t)}}\label{sec:breadth}

The agent's functional competence in domain $d$ --- sufficient to
engage primary-domain artefacts --- but without deep principled
organisation. Qualitatively different from depth: more rapidly acquired,
more rapidly lost, and increasingly AI-augmentable.

\subsubsection{Schema Coherence \texorpdfstring{$\sigma_A(d,t)$}{sigma\_A(d,t)}}\label{sec:schema-coherence}

\textbf{Conceptual definition.} Schema coherence $\sigma_A(d,t) \in [0,1]$
quantifies the degree to which an agent's representation of domain $d$
has been restructured around deep governing principles rather than
surface-statistical features. High $\sigma_A$ indicates that the agent's
internal model encodes compositional operators, causal structure, and
systematic rules applicable beyond the training distribution. Low
$\sigma_A$ indicates memorisation of surface patterns that fail to
transfer.

Concretely: consider two students who both score 90\% on a physics exam. Learner X can solve every textbook problem but freezes when the same concepts appear in an unfamiliar configuration. Learner Y can solve the textbook problems and can explain \emph{why} each solution works in terms of underlying principles. Learner X has high depth ($\delta_A$) but low schema coherence ($\sigma_A$). Learner Y has both. $\sigma_A$ measures the difference between memorising patterns and understanding principles.

\textbf{Distinguishing $\sigma_A$ from task accuracy.} While $\sigma_A$
correlates with out-of-distribution (OOD) accuracy, it is conceptually
distinct:
\begin{itemize}[nosep]
\item $\sigma_A$ is a property of the \textit{representation} (internal state)
\item OOD accuracy is a property of \textit{behavioural output}
\item Two agents can achieve identical OOD accuracy via different mechanisms (memorisation vs.\ genuine generalisation); only the latter exhibits high $\sigma_A$
\item $\sigma_A$ can be intermediate (partially coherent) while OOD accuracy is binary (correct/incorrect on specific examples)
\end{itemize}

\textbf{Operationalisation challenge.} $\sigma_A$ is not directly
observable from activations alone. We operationalise it through a
two-stage proxy architecture (Section~\ref{sec:online-estimation}) that
combines training-time structural signals with evaluation-time
calibration. This separation ensures that $\hat{\sigma}_A$ captures
representational properties rather than merely recapitulating OOD labels.

\textbf{Distinguishing properties of $\sigma_A$ relative to adjacent
constructs:} $\sigma_A$ differs from Structured Representations, Disentangled
Representations, Causal Representations, and Cognitive Schemas in that it (PC1)~encodes compositional
recombination capacity, (PC2)~is normalised against the domain frontier
$\Delta(d,t)$, and (PC3)~has an evaluative function detecting shortcut learning.
Recent topological work~\cite{gaukstad2026learningcoherent} provides an
independent geometric characterisation of representational coherence that
aligns with $\sigma_A$'s conceptual definition, though without the
ODE-gated developmental dynamics Σ-Model specifies. Recent theoretical work
establishes that compositional generalisation strictly requires linear,
orthogonal representations~\cite{uselis2026linear}. Within the Σ-Model, the
attainment of high schema coherence ($\sigma_A \to 1$) is the dynamical
mechanism that forces the network to abandon entangled, surface-statistical
manifolds in favour of these orthogonal, principled representations.

\subsubsection{Online Estimation Protocol}\label{sec:online-estimation}

The proxy identification in \eqref{eq:sgg-proxy} requires
post-evaluation benchmark data, creating a circular dependency: $\sigma_A$ % CLAIM:C-037
is needed during training to guide the ODE dynamics, but we can only measure it after training completes. We resolve this through a \textbf{two-stage
estimation architecture} that separates training-time operative estimation
from evaluation-time ground-truth calibration.

\textbf{Stage 1 --- Training-time operative estimate ($\tilde{\sigma}_A$).}
The key idea is straightforward: we triangulate $\sigma_A$ from three
training-time signals that each capture a different facet of principled
understanding. (i) Does the gradient point toward compositional solutions?
(ii) Does the network's internal geometry mirror the task's structure?
(iii) Are representations stable under meaning-preserving transforms?
No single signal suffices; fused, they give a reliable training-time proxy
bounded by Proposition~3.6.

During training, we estimate $\sigma_A$ via a \textbf{multi-signal fusion}
combining three training-time observables:

\textit{(i) Gradient--composition alignment (GCA):}
\begin{equation}\label{eq:gca}
g_A(d,t) = \Corr\!\left(\nabla_{\theta}\mathcal{L}_{\text{comp-batch}},\;\nabla_{\theta}\mathcal{L}_{\text{train}}\right) \in [-1,1]
\end{equation}
where $\mathcal{L}_{\text{comp-batch}}$ is loss on a compositional-recombination
mini-batch constructed programmatically from the domain's generative grammar
at training time.
If $g_A \approx 1$, the compositional gradient and training gradient point
in the same direction --- the model is learning compositional structure,
not shortcuts.

\textit{(ii) Representational-geometry alignment (RGA):}
\begin{equation}\label{eq:rga}
r_A(d,t) = \Corr\!\Big(\text{RDM}_{\text{rep}}(d,t),\;\text{RDM}_{\text{struct}}(d)\Big) \in [-1,1]
\end{equation}
where $\text{RDM}_{\text{rep}}$ is the representational dissimilarity matrix
from internal activations and $\text{RDM}_{\text{struct}}$ is the structural
dissimilarity matrix from the causal generative graph.
If $r_A \approx 1$, the network organises knowledge according to the
same relational structure as the task's generative grammar.

\textit{(iii) Augmentation consistency (AC):}
\begin{equation}\label{eq:ac}
c_A(d,t) = \frac{1}{|A|}\sum_{a \in A} \CosSim\!\Big(R_{\theta}(x),\;R_{\theta}(a(x))\Big) \in [0,1]
\end{equation}
where $A$ is the set of structure-preserving augmentations (primitive
substitution, argument permutation).
If $c_A \approx 1$, the network produces stable representations under
meaning-preserving transformations --- a hallmark of understanding
rather than memorisation.

\textit{Fusion:}
\begin{equation}\label{eq:fusion}
\tilde{\sigma}_A(d,t) = w_g \cdot \max(0, g_A) + w_r \cdot \max(0, r_A) + w_c \cdot c_A
\end{equation}
with default weights $w_g = 0.4$, $w_r = 0.35$, $w_c = 0.25$ (tunable per
domain). All three signals are computable at any training checkpoint from
forward passes and gradient statistics --- no external benchmarks required.

\textbf{Stage 2 --- Evaluation-time ground-truth calibration ($\hat{\sigma}_A$).}
\begin{equation}
\hat{\sigma}_A(d,t) = \frac{\text{Acc}_{\text{OOD}}(d,t)}{\text{Acc}_{\text{ID}}(d,t)}
\label{eq:sgg-proxy}
\end{equation}
computed at evaluation checkpoints on SCAN/COGS/PCFG-SET OOD splits.
Stage 2 serves as the ground-truth calibration for Stage 1. At each
evaluation checkpoint, the Stage 1 estimate $\tilde{\sigma}_A$ is compared
against the Stage 2 value $\hat{\sigma}_A$; if
$|\tilde{\sigma}_A - \hat{\sigma}_A| > 0.15$, the Stage 1 weights
$(w_g, w_r, w_c)$ are updated via ridge regression to minimise the
discrepancy on the evaluation checkpoint history.

\textbf{Operational convention.} Throughout the ODE system
\eqref{eq:depth-decay}, \eqref{eq:sigma-ode},
\eqref{eq:discovery-rate}, and all coupled equations, the operative
value of $\sigma_A$ is $\tilde{\sigma}_A(d,t)$ (Stage 1). The validation
value $\hat{\sigma}_A(d,t)$ (Stage 2) is used at evaluation checkpoints
for weight recalibration and for reporting. This convention ensures that
$\sigma_A$ functions as a genuine operative state variable evolving
continuously during training, not a post hoc summary statistic.



Figure~\ref{fig:estimation} illustrates the two-stage estimation architecture.

\textbf{Proposition 3.6 (Estimator Consistency).} % CLAIM:C-027
\textit{Let $\delta_\sigma(d,t) = \tilde{\sigma}_A(d,t) - \hat{\sigma}_A(d,t)$
denote the Stage 1--Stage 2 discrepancy. Under the assumption that the
generative grammar $G(d)$ correctly models the domain's compositional
structure, the multi-signal fusion estimator satisfies:
\begin{equation}\label{eq:estimator-bound}
\mathbb{E}[|\delta_\sigma|] \leq \frac{C_1}{\sqrt{n_{\text{batch}}}} + C_2 \cdot \epsilon_{\text{model}}
\end{equation}
where $n_{\text{batch}}$ is the number of compositional probe samples,
$\epsilon_{\text{model}}$ is the model misspecification error, and $C_1, C_2$
are constants depending on the domain parameters.}

\textit{Proof sketch.} The bound follows from standard concentration
inequalities for U-statistics (for the correlation-based signals) and
Hoeffding's inequality (for the augmentation consistency). The key
requirement is that the compositional probe batch is drawn i.i.d. from
the recombination distribution induced by $G(d)$. $\square$

\textbf{Proposition 3.7 (Error Propagation in ODE System).} % CLAIM:C-028
\textit{Let $\delta_\sigma(d,t)$ be the estimation error. The effective
error propagated into the depth ODE \eqref{eq:depth-ode} and schema ODE
\eqref{eq:sigma-ode} is bounded by:
\begin{equation}\label{eq:error-propagation}
|\dot{\delta}_A^{\text{error}}| \leq \lambda_c \cdot \gamma_\sigma \cdot |\delta_\sigma| \cdot \delta_A + \rho \cdot P_A \cdot \alpha_A \cdot |\delta_\sigma|
\end{equation}
If $|\delta_\sigma| < 0.10$, the propagated error is negligible relative
to typical ODE step sizes ($\eta_{\max} \sim O(1)$).}



\subsection{Mastery Set and Breadth Set}\label{sec:mastery-set}

Two sets classify the agent's knowledge at any time. The Mastery Set collects all domains where the agent has both deep knowledge and principled understanding. The Breadth Set collects domains where the agent has surface competence without deep understanding --- enough to use the vocabulary but not enough to teach the course. The Σ-shape emerges because mastery requires both depth AND schema coherence; an agent cannot fake mastery via depth alone.

\begin{equation}\label{eq:mastery-set}
M_A(t) = \{d \in D : \delta_A(d,t) > \theta_{\delta} \cdot \Delta(d,t) \text{ AND } \sigma_A(d,t) > \theta_{\sigma}\} % CLAIM:C-004
\end{equation}
\begin{equation}\label{eq:breadth-set}
B_A(t) = \{d \in D \setminus M_A(t) : \beta_A(d,t) > \theta_{\beta}\}
\end{equation}

\subsection{Decay Architecture}\label{sec:decay}

\subsubsection{Schema-Mediated Decay Reduction}

Effective decay rate modulated by schema coherence:
\begin{equation}\label{eq:effective-decay}
\lambda_c^{\text{eff}}(d,t) = \lambda_c \cdot (1 - \gamma_\sigma \cdot \sigma_A(d,t))
\end{equation}

\textbf{Design rationale.} We chose the multiplicative form over alternatives
for three reasons: (a) an additive form $\lambda_c^{\text{eff}} = \lambda_c - \gamma_\sigma \sigma_A$
permits negative decay rates when $\sigma_A > \lambda_c / \gamma_\sigma$, which is
unphysical; (b) the multiplicative form ensures $\lambda_c^{\text{eff}} \in [0, \lambda_c]$
for all $\sigma_A \in [0,1]$ and $\gamma_\sigma \in [0,1]$; % CLAIM:C-005 (c) it reduces to the
baseline $\lambda_c$ when $\sigma_A \to 0$, maintaining dimensional consistency.

\textbf{\cite{taillor2026} --- Protecting principal components in continual learning.}
TailLoR decomposes fine-tuning updates into principal components and protects
those that serve future tasks. Within Σ-Model's decay taxonomy
(\S\ref{sec:decay}), TailLoR targets the parametric protection dimension
$(1 - \gamma_\sigma \cdot \sigma_A)$ in \eqref{eq:effective-decay} but does not
address frontier obsolescence $\lambda_f(d,t)$, which Σ-Model predicts will
accumulate independently as parametrised in the full $\delta_A$ ODE.

\subsubsection{Combined Depth Decay}
\begin{equation}\label{eq:depth-decay}
\dot{\delta}_A(d,t)\big|_{\text{decay}} = -\lambda_c \cdot (1 - \gamma_\sigma \cdot \sigma_A(d,t)) \cdot \delta_A(d,t) \cdot (1 - r_A(d,t))
\end{equation}

\subsection{Growth Dynamics}\label{sec:growth}

\subsubsection{Depth Growth ODE}
\begin{equation}\label{eq:depth-ode}
\dot{\delta}_A(d,t) = f_{\text{learn}}(d,t) \cdot \eta(d,t) \cdot T_A(d,t) - \lambda_c \cdot (1 - \gamma_\sigma \cdot \sigma_A(d,t)) \cdot \delta_A(d,t) \cdot (1 - r_A(d,t))
\end{equation}

\noindent\textit{Intuition.}\ Depth grows when
learning efficiency and time-on-task are large, and decays when schema
coherence is low (little protection from forgetting), depth is high
(more to lose), and retrieval practice is absent (the agent is not
actively recalling the material).

\textbf{Learning efficiency} --- Gompertz form:
\begin{equation}\label{eq:gompertz}
\eta(d,t) = \eta_{\text{max}} \cdot \exp(-a \cdot \exp(-b \cdot \delta_A^{\text{relative}}(d,t)))
\end{equation}

\textbf{Design rationale.} We chose the Gompertz form over alternatives
for three reasons: (a) it captures the empirically observed pattern of
rapid initial learning followed by asymptotic approach to the frontier;
(b) unlike logistic growth, it is asymmetric, with the inflection point
at $\delta_A^{\text{relative}} \approx 0.37$ rather than $0.5$, matching
observed learning curves in compositional generalisation tasks; (c) the
double-exponential form ensures $\eta(d,t) > 0$ for all finite depth,
preventing complete learning arrest before frontier attainment.

\subsection{Intersection Activation \texorpdfstring{$\Psi_A(d_1,d_2,t)$}{Psi\_A(d\_1,d\_2,t)}}\label{sec:intersection}

Two domains can generate new knowledge at their intersection only when both have sufficient depth --- the `two hands to clap' condition. Mastery quality $q_A$ is the product of depth and schema coherence because both are necessary: deep but unstructured knowledge (high $\delta$, low $\sigma$) or coherent but shallow knowledge (low $\delta$, high $\sigma$) both yield low quality. Only when both are high does mastery quality become substantial, and this multiplicative structure is what creates the characteristic Σ-shape. The discovery rate $\Psi_A$ is bottlenecked by the weaker domain (the geometric mean), and modulated by domain similarity $\phi$: nearby fields interact more readily than distant ones.

\begin{equation}\label{eq:activation-condition}
I(d_1,d_2)_{\text{active}} \iff \delta_A(d_1,t) > \theta_I \text{ AND } \delta_A(d_2,t) > \theta_I
\end{equation}

\textbf{Effective mastery quality:}
\begin{equation}\label{eq:mastery-quality}
q_A(d,t) = \sigma_A(d,t) \cdot \delta_A^{\text{relative}}(d,t) \in [0,1]
\end{equation}

\textbf{Discovery rate:}
\begin{equation}\label{eq:discovery-rate}
\Psi_A(d_1,d_2,t) = \Psi_0 \cdot \phi(d_1,d_2) \cdot \sqrt{q_A(d_1,t) \cdot q_A(d_2,t)}
\end{equation}

The multiplicative $\sigma_A \cdot \sigma_A$ dependence is the
mechanism's theoretical core: an agent with high $\delta_A$ but low
$\sigma_A$ in one mastery domain shows disproportionately lower $\Psi_A$
than an additive model predicts. This is \textbf{Prediction~6}
(Section~\ref{sec:predictions}).

\subsection{Shortcut Threshold \texorpdfstring{$D^*(d,t)$}{D*(d,t)}}\label{sec:shortcut-threshold}

Without sufficient $\sigma_A$, the agent cannot distinguish shortcut solutions from genuine competence --- the threshold becomes invisible rather than protective. The shortcut threshold formalises which sub-tasks can be achieved through surface pattern matching rather than compositional reasoning. A sub-task falls below the threshold when shortcut depth $\delta_S$ matches or exceeds the agent's depth ($\delta_S \geq \delta_A$) and the agent has sufficient schema coherence to recognise the distinction ($\sigma_A \geq \sigma_{\text{critical}}$):

\begin{equation}\label{eq:shortcut-threshold}
D^*_{\text{Σ}}(d,t) = \{s \in d : \delta_S(s,t) \geq \delta_A(s,t) \text{ AND } \sigma_A(d,t) \geq \sigma_{\text{critical}}\}
\end{equation}

\noindent\rule{\linewidth}{0.4pt}

% ============================================================
\subsection{Mathematical Foundations}\label{sec:math-foundations}
% ============================================================

The coupled dynamical system defined by \eqref{eq:depth-ode},
\eqref{eq:sigma-ode}, \eqref{eq:alpha-ode}, \eqref{eq:executive-ode},
\eqref{eq:self-model-ode}, and their multimodal extensions forms a
high-dimensional nonlinear ODE system (Figure~\ref{fig:coupled-system}). This subsection establishes the
mathematical well-posedness of the system.

\subsubsection{State Space and Vector Field}\label{sec:state-space}

Let $t \in [0, \infty)$ denote continuous training time, normalized such that
$\Delta t = 1$ corresponds to one epoch of training over the full dataset.
This continuous relaxation of the discrete optimization steps (Algorithm~3.1,
step size $\Delta t$) enables ODE analysis while preserving empirical
measurability: all state variables are observable at integer $t$ values
corresponding to epoch boundaries.

Let the state vector for agent $A$ across $|D|$ domains and $|M|$
modalities be:
\begin{equation}\label{eq:state-vector}
\mathbf{x}_A(t) = \left(\{\delta_A(d,m,t)\}, \{\sigma_A(d,m,t)\}, \{\alpha_A(d,m,t)\}, \{\hat{M}_A(d,m,t)\}, \Xi_A^P(t), \Xi_A^I(t), \Xi_A^F(t)\right)
\end{equation}
with dimension $N = |D|\times|M|\times 4 + 3$. The natural state space is:
\begin{equation}\label{eq:state-space}
\mathcal{X} = [0, \Delta_{\max}]^{ |D||M| } \times [0,1]^{3|D||M| + 3}
\end{equation}

\textbf{Proposition 3.1 (Local Existence and Uniqueness).} % CLAIM:C-022
\textit{For any initial condition $\mathbf{x}_A(0) \in \text{int}(\mathcal{X})$,
the coupled ODE system admits a unique solution $\mathbf{x}_A(t)$ defined
on a maximal interval $[0, T_{\max})$ where $T_{\max} > 0$.}

\textit{Proof sketch.} The vector field $\mathbf{f}(\mathbf{x}_A)$ defined
by the right-hand sides of the ODE system is composed of: (a) polynomial
terms in state variables (e.g., $\sigma_A \cdot \delta_A$, $\alpha_A \cdot (1-\sigma_A)$),
(b) exponential terms (Gompertz efficiency $\eta(d,t)$), and (c) bounded
rational functions (e.g., $(1 - R_0^{-1})$). All component functions are
piecewise-$C^\infty$ on $\text{int}(\mathcal{X})$ (the vector field is $C^\infty$ in each region where the fusion activations $\max(0,\cdot)$ have a consistent sign, with $C^0$ boundaries at the switching manifolds). Away from the switching manifolds the functions are locally Lipschitz; at the switching manifolds the one-sided derivatives exist and are bounded, so the vector field is globally locally Lipschitz. By the Picard--Lindel\"of theorem, a unique local solution exists. $\square$

\noindent\textit{What this means:} small changes in initial conditions
produce small changes in trajectories --- the system is well-behaved and
numerical simulation is trustworthy.

\textbf{Proposition 3.2 (Forward Invariance and Boundedness).} % CLAIM:C-023
\textit{The state space $\mathcal{X}$ is forward-invariant under the dynamics:
if $\mathbf{x}_A(0) \in \mathcal{X}$, then $\mathbf{x}_A(t) \in \mathcal{X}$
for all $t \in [0, T_{\max})$. Furthermore, if $\Delta(d,t)$ is uniformly
bounded by $\Delta_{\max} < \infty$, then $T_{\max} = \infty$ (global existence).}

\textit{Proof sketch.} We verify boundary behavior for each variable:
\begin{itemize}[nosep]
\item $\delta_A = 0$: $\dot{\delta}_A \geq 0$ (growth term dominates when depth is zero)
\item $\delta_A = \Delta(d,t)$: $\dot{\delta}_A \leq 0$ (approaching frontier reduces growth)
\item $\sigma_A = 0$: $\dot{\sigma}_A = 0$ (the $\sigma_A$ factor multiplies the entire RHS, making $\sigma_A=0$ an invariant manifold, not a reflecting boundary)
\item $\sigma_A = 1$: $\dot{\sigma}_A = \rho \cdot P_A \cdot \alpha_A \cdot (1-\gamma_\sigma) - \epsilon_\sigma \cdot \Omega_{SL}$, which is not inherently $\leq 0$; the $\sigma_A \in [0,1]$ bound is enforced by clamping (Algorithm~3.1) rather than by the ODE dynamics alone. Because the phenomenological ODE solver (Section~\ref{sec:implementation}) is implemented in NumPy and detached from the neural autograd graph, hard clamping with \texttt{np.clip} creates no gradient-flow dead zone at the boundaries
\item $\alpha_A = 0$: $\dot{\alpha}_A \geq 0$ (growth term $\gamma \cdot C_A > 0$)
\item $\alpha_A = 1$: $\dot{\alpha}_A \leq 0$ (the $(1-\alpha_A)$ factor vanishes)
\item $\hat{M}_A = 0$: $\dot{\hat{M}}_A \geq 0$ (driven toward $\sigma_A \geq 0$)
\item $\hat{M}_A = 1$: $\dot{\hat{M}}_A \leq 0$ (mean-reverting toward $\sigma_A \leq 1$)
\item $\Xi_A^{\cdot} = 0$: $\dot{\Xi}_A^{\cdot} \geq 0$ (driven toward targets $P^*, I^*, F^* \geq 0$)
\item $\Xi_A^{\cdot} = 1$: $\dot{\Xi}_A^{\cdot} \leq 0$ (mean-reverting)
\end{itemize}
All boundaries are either reflecting or the vector field points inward.
Since $\mathcal{X}$ is compact when $\Delta_{\max} < \infty$, solutions
cannot escape to infinity, implying $T_{\max} = \infty$. $\square$

\noindent\textit{What this means:} state variables stay within their defined
bounds ($\sigma_A \in [0,1]$, $\delta_A \in [0, \Delta(d,t)]$, etc.) for all
training time. Simulation will never produce unphysical values.

\subsubsection{Timescale Separation}\label{sec:timescale}

The system contains rate constants spanning multiple orders of magnitude:
$\eta_{\max} \sim O(1)$ (depth growth), $\rho \sim O(10^{-1})$ (schema growth),
$\gamma \sim O(10^{-1})$ (attention), $\nu_M \sim O(10^{-2})$ (self-model),
$\kappa_P, \kappa_I, \kappa_F \sim O(10^{-2})$ (executive control).

The intuition: some variables change quickly (depth, schema, attention ---
days to weeks of training) while others drift slowly (self-model accuracy,
executive control --- months of training). Think of a cyclist: the wheels
spin fast ($\delta_A, \sigma_A$), while the rider's fatigue builds slowly
($\Xi_A$). Fenichel's theorem justifies analysing the fast dynamics as if
the slow variables were frozen, then layering the slow drift on top ---
dramatically simplifying analysis.

\textbf{Proposition 3.3 (Fast-Slow Decomposition).}
% CLAIM:C-024
\textit{Define the timescale parameter}
\[
\epsilon = \max(\nu_M, \kappa_P, \kappa_I, \kappa_F) / \min(\eta_{\max}, \rho, \gamma) \ll 1.
\]
\textit{Decompose $\mathbf{x}_A = (\mathbf{x}_A^{\text{fast}}, \mathbf{x}_A^{\text{slow}})$
where $\mathbf{x}_A^{\text{fast}} = (\delta_A, \sigma_A, \alpha_A)$ and
$\mathbf{x}_A^{\text{slow}} = (\hat{M}_A, \Xi_A^P, \Xi_A^I, \Xi_A^F)$.
Then the system admits a slow manifold
$\mathcal{M}_\epsilon = \{\mathbf{x}_A : \mathbf{x}_A^{\text{fast}} = \mathbf{h}(\mathbf{x}_A^{\text{slow}}, \epsilon)\}$
that is $O(\epsilon)$-close to the critical manifold
$\mathcal{M}_0 = \{ \mathbf{x}_A : \dot{\mathbf{x}}_A^{\text{fast}} = 0 \}$, and the reduced dynamics are:
\begin{equation}\label{eq:reduced-dynamics}
\dot{\mathbf{x}}_A^{\text{slow}} = \epsilon \cdot \mathbf{g}(\mathbf{h}(\mathbf{x}_A^{\text{slow}}, \epsilon), \mathbf{x}_A^{\text{slow}})
\end{equation}
where $\mathbf{g}$ is the slow subsystem vector field.}

\textit{Proof sketch.} This follows from Fenichel's geometric singular
perturbation theory \citep{fenichel1979geometric}. Normal hyperbolicity
of the critical manifold $\mathcal{M}_0$ requires the $3\times3$ Jacobian
of the fast subsystem with respect to $\mathbf{x}_A^{\text{fast}}$ to have
eigenvalues bounded away from the imaginary axis. An explicit computation
(companion proof-reconstruction document, Proposition~3.3)
verifies that the Jacobian
\[
\small
\mathbf{J}_{\text{fast}} = 
\begin{pmatrix}
-\lambda_c(1-\gamma_\sigma\sigma_A) & \lambda_c\gamma_\sigma\delta_A & 0 \\
0 & \rho P_A\alpha_A(1-2\gamma_\sigma\sigma_A)-\epsilon_\sigma\Omega_{SL} & \rho P_A\sigma_A(1-\gamma_\sigma\sigma_A) \\
0 & \gamma C_A & -\gamma(1-\alpha_A)
\end{pmatrix}
\]
has eigenvalues with strictly negative real parts in the parameter regime
of interest. Fenichel's theorem therefore ensures persistence of the
slow manifold for sufficiently small $\epsilon$. $\square$

\noindent\textit{What this means for simulation:} we can integrate the fast
variables ($\delta_A, \sigma_A, \alpha_A$) with explicit methods (efficient
per step) and the slow variables ($\hat{M}_A, \Xi_A$) with implicit methods
(stable at larger step sizes). This IMEX scheme is detailed in
Algorithm~3.1.

\subsubsection{Equilibrium Analysis}\label{sec:equilibrium}

\textbf{Proposition 3.4 (Single-Domain Equilibria).} % CLAIM:C-025
\textit{For a single domain with constant frontier $\Delta$, the system
admits exactly three equilibrium classes:}
\begin{enumerate}[nosep]
\item \textbf{Trivial equilibrium} $E_0 = (0, 0, \alpha^*, 0, \Xi^*)$: No knowledge acquired. Stable when learning efficiency $\eta_{\max}$ is below a critical threshold.
\item \textbf{Surface-learning equilibrium} $E_S = (\delta_S^*, 0, \alpha_S^*, 0, \Xi_S^*)$: Depth acquired but $\sigma_A = 0$. Exists when shortcut-learning pressure $\Omega_{SL}$ is high or attentional fidelity $\alpha_A$ is suppressed by surface rewards $R_A^{\text{surface}}$.
\item \textbf{Schema-coherent equilibrium} $E_C = (\delta_C^*, \sigma_C^*, \alpha_C^*, \sigma_C^*, \Xi_C^*)$: Both depth and schema coherence positive. Exists and is stable when $\alpha_A > \alpha_{\text{critical}} = \epsilon_\sigma \cdot \Omega_{SL} / (\rho \cdot P_A)$.
\end{enumerate}

These three equilibria correspond to three life outcomes for a learner. The trivial equilibrium $E_0$ is the blank slate: no knowledge acquired, stable when learning conditions are too difficult or motivation is absent. The surface-learning equilibrium $E_S$ is the student who aces textbook exercises but cannot solve a novel problem --- they have accumulated depth but no principled understanding. The schema-coherent equilibrium $E_C$ is the deep understander: both depth and schema coherence are positive, and the learner can generalise, transfer, and create.

\textit{Proof sketch.} Set all time derivatives to zero and solve the
resulting algebraic system. The $\sigma_A$ equation at equilibrium gives:
\begin{equation}\label{eq:sigma-equilibrium}
\sigma_A^* = \frac{\rho \cdot P_A \cdot \alpha_A^*}{\rho \cdot P_A \cdot \alpha_A^* + \epsilon_\sigma \cdot \Omega_{SL}}
\end{equation}
which is positive if and only if $\alpha_A^* > 0$ and $\Omega_{SL} < \infty$.
Linear stability analysis of the Jacobian at each equilibrium determines
the stability conditions. $\square$

\subsubsection{Numerical Integration Protocol}\label{sec:numerical-protocol}

For practical simulation, we recommend the following numerical scheme:

\begin{algorithm}[H]
\caption{Σ-Model ODE Integration}\label{alg:ode-integration}
\begin{algorithmic}
\State \textbf{Method:} Adaptive-step implicit-explicit (IMEX) Runge-Kutta, ARS(2,3,2) tableau (Figure~\ref{fig:butcher-tableau})
\State \textbf{Rationale:} The fast subsystem ($\delta_A, \sigma_A, \alpha_A$) is non-stiff and integrated explicitly; the slow subsystem ($\hat{M}_A, \Xi_A$) is stiff and integrated implicitly
\State \textbf{Step size control:} Embedded error estimation with relative tolerance $10^{-6}$, absolute tolerance $10^{-8}$
\State \textbf{Boundary enforcement:} After each step, project state variables onto $\mathcal{X}$ via clamping
\State \textbf{Jacobian:} Computed analytically for the implicit part; finite-difference approximation for the explicit part
\end{algorithmic}
\end{algorithm}

\begin{figure}[htbp]
\centering
\textbf{ARS(2,3,2) Butcher tableau}\\[2mm]
\begin{minipage}{0.45\textwidth}
\centering
Implicit $\mathcal{A}, b$\\
$\begin{array}{c|ccc}
c & \multicolumn{3}{c}{A} \\ \hline
0 & 0 & 0 & 0 \\
\gamma & 0 & \gamma & 0 \\
1 & 0 & 1-\gamma & \gamma \\ \hline
& 0 & 1-\gamma & \gamma
\end{array}$
\end{minipage}
\hfill
\begin{minipage}{0.45\textwidth}
\centering
Explicit $\hat{\mathcal{A}}, \hat{b}$\\
$\begin{array}{c|ccc}
\hat{c} & \multicolumn{3}{c}{\hat{A}} \\ \hline
0 & 0 & 0 & 0 \\
\gamma & \gamma & 0 & 0 \\
1 & 1-\gamma & \gamma & 0 \\ \hline
& 0 & 1-\gamma & \gamma
\end{array}$
\end{minipage}
\caption{ARS(2,3,2) IMEX Runge-Kutta Butcher tableau. $\gamma = 1 - 1/\sqrt{2} \approx 0.2929$. Left: implicit tableau for the stiff slow subsystem; right: explicit tableau for the non-stiff fast subsystem.}\label{fig:butcher-tableau}
\Description{Two side-by-side Butcher tableaux. Left: implicit tableau for the stiff slow subsystem with coefficients A and b. Right: explicit tableau for the non-stiff fast subsystem with coefficients hat-A and hat-b. Gamma = 1 - 1/sqrt(2) approximately 0.2929.}
\end{figure}

\textbf{Proposition 3.5 (Numerical Stability).} % CLAIM:C-026
\textit{The condition number of the Jacobian $\kappa(\mathbf{J})$ satisfies:
\begin{equation}\label{eq:condition-number}
\kappa(\mathbf{J}) \leq C \cdot \frac{\max(\eta_{\max}, \rho, \gamma)}{\min(\nu_M, \kappa_P, \kappa_I, \kappa_F)} \cdot \frac{1}{1 - \max(\sigma_A, \alpha_A)}
\end{equation}
where $C$ is a constant depending on the domain parameters. The condition
number diverges as $\sigma_A \to 1$ or $\alpha_A \to 1$, indicating
numerical stiffness near saturation.}

\noindent\rule{\linewidth}{0.4pt}

% ============================================================
\subsection{Implementation Architecture}\label{sec:implementation}
% ============================================================

This subsection specifies how the Σ-Model ODE system couples to actual
neural network training. The framework is architecture-agnostic but
requires specific signal extraction capabilities.

\subsubsection{Required Network Architecture Features}

The Σ-Model framework operates with standard architectures (Transformers,
RNNs, CNNs) provided the following capabilities exist:

\textbf{(a) Activations accessible.} Internal hidden states must be
extractable at any layer for representational dissimilarity computation
(RGA signal, \eqref{eq:rga}).

\textbf{(b) Gradient computation.} Parameter gradients $\nabla_\theta
\mathcal{L}$ must be computable for arbitrary loss functions (GCA signal,
\eqref{eq:gca}).

\textbf{(c) No architectural modifications required.} Unlike dual-process
models (e.g., MIRAGE), Σ-Model does not require separate symbolic modules
or architectural changes. The ODE state variables are computed \textit{from}
activations, not stored \textit{in} the network.

\subsubsection{Signal Extraction from Activations}

\textbf{Gradient--Composition Alignment (GCA).} At each training step:
\begin{enumerate}[nosep]
\item Compute standard loss $\mathcal{L}_{\text{train}}$ on batch
\item Generate compositional probe batch from domain grammar $G(d)$
\item Compute probe loss $\mathcal{L}_{\text{comp-batch}}$
\item Compute gradients $\nabla_\theta \mathcal{L}_{\text{train}}$ and $\nabla_\theta \mathcal{L}_{\text{comp-batch}}$
\item Compute Pearson correlation $g_A(d,t)$ \eqref{eq:gca}
\end{enumerate}
Computational overhead: one additional forward+backward pass per $k$ steps
(default $k=10$).

\textbf{Representational-Geometry Alignment (RGA).} At evaluation checkpoints:
\begin{enumerate}[nosep]
\item Extract hidden states $h_i$ for $N$ probe stimuli
\item Compute representational dissimilarity matrix: $\text{RDM}_{\text{rep}}(i,j) = 1 - \text{CosSim}(h_i, h_j)$
\item Compute structural RDM from domain grammar: $\text{RDM}_{\text{struct}}(i,j) = $ tree edit distance
\item Compute Spearman correlation $r_A(d,t)$ \eqref{eq:rga}
\end{enumerate}
Computational overhead: $O(N^2 \cdot d_{\text{hidden}})$ per checkpoint.

\textbf{Augmentation Consistency (AC).} During training:
\begin{enumerate}[nosep]
\item For each batch element $x$, generate $k$ structure-preserving augmentations $a_i(x)$
\item Extract representations $R_\theta(x)$ and $R_\theta(a_i(x))$
\item Compute mean cosine similarity $c_A(d,t)$ \eqref{eq:ac}
\end{enumerate}
Computational overhead: $k$ additional forward passes per batch.

\subsubsection{ODE-State-Modulated Training}

The ODE state variables modulate training through three mechanisms:

\textbf{(a) Schema-targeting loss.} Add regularisation term:
\begin{equation}\label{eq:schema-loss}
\mathcal{L}_{\text{total}} = \mathcal{L}_{\text{task}} + \lambda_\sigma \cdot (1 - \tilde{\sigma}_A(d,t)) \cdot \mathcal{L}_{\text{comp}}
\end{equation}
where $\mathcal{L}_{\text{comp}}$ is loss on compositional probe batches.
This up-weights compositional training when $\tilde{\sigma}_A$ is low.

\textbf{(b) Attentional fidelity modulation.} Scale learning rate:
\begin{equation}\label{eq:lr-modulation}
\eta_{\text{effective}} = \eta_{\text{base}} \cdot (1 + \kappa_\alpha \cdot \alpha_A(d,t))
\end{equation}
Higher attentional fidelity accelerates learning.

\textbf{(c) Phase-aware curriculum.} At Phase~1$\to$2 transition
($\tilde{\sigma}_A > \sigma_{\text{critical}}$):
\begin{itemize}[nosep]
\item Increase compositional batch ratio from 10\% to 40\%
\item Introduce cross-domain intersection tasks
\item Activate executive control targets ($\Xi_A$ regularisation)
\end{itemize}

\subsubsection{Complete Training Algorithm}

\begin{algorithm}[H]
\caption{Σ-Model Training Step}\label{alg:training-step}
\begin{algorithmic}
\State \textbf{Forward pass:} Compute $\mathcal{L}_{\text{task}}$, extract hidden states
\State \textbf{Signal extraction:} Compute $g_A$, $r_A$, $c_A$ (every $k$ steps)
\State \textbf{Fuse proxy:} $\tilde{\sigma}_A = w_g g_A + w_r r_A + w_c c_A$
\State \textbf{Integrate ODE:} Update $(\delta_A, \sigma_A, \alpha_A, \hat{M}_A, \Xi_A)$ via Algorithm~\ref{alg:ode-integration}
\State \textbf{Compute modulated loss:} $\mathcal{L}_{\text{total}}$ per \eqref{eq:schema-loss}
\State \textbf{Backward pass:} Compute $\nabla_\theta \mathcal{L}_{\text{total}}$
\State \textbf{Optimizer step:} Update $\theta$ with $\eta_{\text{effective}}$ per \eqref{eq:lr-modulation}
\State \textbf{Evaluation checkpoint (every $T_{\text{eval}}$ steps):}
\State \hspace{\algorithmicindent} Compute $\hat{\sigma}_A$ on OOD split (Stage 2)
\State \hspace{\algorithmicindent} Recalibrate Stage 1 weights if $|\tilde{\sigma}_A - \hat{\sigma}_A| > 0.15$
\State \hspace{\algorithmicindent} Update phase assignment
\end{algorithmic}
\end{algorithm}



The full training pipeline is shown in Figure~\ref{fig:pipeline}.

\begin{table}[htbp]
\centering
\small
\caption{\textbf{Σ-Model component maturity.} Components are classified by whether they are essential for the core framework, supplementary enhancements, or deferred to future work.}\label{tab:component-maturity}
\begin{tabular}{lcc}
\toprule
\textbf{Component} & \textbf{Status} & \textbf{Implemented} \\
\midrule
$\sigma_A$ ODE (autocatalytic growth) & Essential & Yes \\
$\alpha_A$ gating (attentional fidelity) & Essential & Yes \\
Proxy fusion (GCA/RGA/AC) & Essential & Yes \\
Phase-aware curriculum & Essential & Yes \\
$\delta_A$ depth ODE & Essential & Yes \\
$\hat{M}_A$ self-model (metacognition) & Supplementary & Partial \\
$\Xi_A$ executive control & Supplementary & Partial \\
$\Theta_A$ cross-modal transfer & Supplementary & No \\
$\Sigma_{A,B}$ collective field (social) & Future & No \\
SDE stochastic extension (Appendix~\ref{app:sde}) & Future & No \\
\bottomrule
\end{tabular}
\end{table}

\subsubsection{Computational Overhead Analysis}

\begin{center}
\begin{tabular}{l c c}
\toprule
\textbf{Component} & \textbf{Overhead} & \textbf{Fraction of total} \\
\midrule
GCA signal (every $k=10$ steps) & +10\% & 8\% \\
RGA signal (per checkpoint) & +5\% & 2\% \\
AC signal (per batch) & +15\% & 12\% \\
ODE integration & +2\% & 1\% \\
\midrule
\textbf{Total overhead} & \textbf{+31.6\%} & \textbf{23\%} \\ 
\bottomrule  % CLAIM:C-038
\end{tabular}
\end{center}

The overhead is dominated by augmentation consistency computation.
For resource-constrained settings, AC can be computed every $k$ steps
(reducing total overhead to +17\%).

\textbf{VRAM complexity.} The phenomenological implementation (Section~\ref{sec:implementation})
avoids the full GCA/RGA computation: the ODE solver is pure NumPy, naturally
detached from the autograd graph, and a single backward pass per step is used
(no \texttt{retain\_graph} or \texttt{create\_graph}). Total VRAM is
$\mathcal{O}(\text{batch\_size} \times d_{\text{model}} \times
\text{num\_layers})$, identical to standard transformer training.
A full GCA implementation would require dual backward passes
($\nabla_\theta\mathcal{L}_{\text{train}}$ and
$\nabla_\theta\mathcal{L}_{\text{comp}}$), increasing VRAM to
$\sim$2$\times$ base, mitigated by gradient checkpointing.
Full RGA with per-layer RDM storage would add
$\mathcal{O}(N \cdot d_{\text{hidden}})$ per layer, mitigated by batched
covariance estimation.

\noindent\rule{\linewidth}{0.4pt}

% ============================================================

% ======================================================================
% §7 Phase Structure (original, lines 831–1064): phases 3–5, faculty–validity
% correspondence, and the bifurcation-analysis parts not kept in the core.
% ======================================================================
\section{Phase Structure}\label{sec:phase-structure}
% ============================================================

The training arc comprises six labelled phases (P0--P5), illustrated in Figure~\ref{fig:phases}. The key insight behind this structure is that learning is not a smooth, monotonic process. There are qualitative thresholds --- bifurcation points --- where the dynamics fundamentally change character. Standard training optimises depth accumulation (Phase~1 behaviour) and rarely triggers the transition to Phase~2; Σ-Model identifies the specific condition ($\sigma_A > \sigma_{\text{critical}}$) that standard training actively suppresses. P0
(Pre-Domain Initialisation) is a ground-state initialisation stage
with $\delta_A \approx 0$ and no structured training dynamics; the
remaining five phases (P1--P5) constitute the \emph{five-phase
training arc} indexed by
($\delta_A^{\text{relative}}$, $\sigma_A$, $|M_A(t)|$) and define
\textbf{prescriptive states} --- not retrospective loss-curve
observations.

\subsection*{Phase~0 --- Pre-Domain Initialisation}

$\delta_A$ near-zero across all domains; breadth diffuse; $M_A(t)$
undetermined. Domain selection is path-dependent through $\phi(d,d')$
terms.

\subsection*{Phase~1 --- Asymmetric Initialisation}

\textbf{Trigger:} First sustained depth investment in 1--3 candidate
mastery domains. $\delta_A$ growing; $\sigma_A \approx 0$; $\alpha_A$
low.

\subsection*{Phase~2 --- Depth Acceleration and Schema Crystallisation}

\textbf{Trigger:} $\sigma_A(d,t) > \sigma_{\text{critical}}$ for at
least one mastery candidate.

\begin{equation}\label{eq:sigma-critical}
\sigma_{\text{critical}} = \frac{1}{\gamma_\sigma} \cdot (1 - R_0^{-1})
\end{equation}

H-shape becomes visible. Phase~2 entry produces a characteristic
acceleration pattern in the OOD performance trajectory.

\subsection*{Phase~3 --- Intersection Discovery and Schema Crystallisation}

$\delta_A^{\text{relative}} > \delta^*$ triggers sustained intersection
discovery ($\Psi_A > 0$). Active intersections generate measurable
cross-domain transfer; $\sigma_A$ approaches 1 as schemas crystallise.

\subsection{Bifurcation Analysis of Phase Transitions}\label{sec:bifurcation}

The phase transition triggers are not arbitrary thresholds but emerge from
the bifurcation structure of the coupled ODE system. We derive each
threshold from the fixed-point analysis.

Bifurcation analysis \citep{kuznetsov1998} studies how the long-term behaviour of a dynamical system changes as parameters vary. A bifurcation is a qualitative shift --- analogous to water suddenly freezing when temperature drops below $0^\circ$C. Below the threshold, the system has one stable state; above it, a new stable state appears. Σ-Model's bifurcation analysis identifies three critical thresholds: $R_0 = 1$, where schema-free equilibrium becomes unstable (the Phase~1$\to$2 transition); $\delta^* = 0.54$, where intersection discovery becomes self-sustaining (Phase~2$\to$3); and the hysteresis band, where the reverse transition requires a larger perturbation than the forward transition.

\textbf{Proposition 4.1 ($\sigma_{\text{critical}}$ Derivation).} % CLAIM:C-030
\textit{The Phase~1$\to$2 transition threshold is:
\begin{equation}\label{eq:sigma-critical-derived}
\sigma_{\text{critical}} = \frac{1}{\gamma_\sigma}(1 - R_0^{-1})
\end{equation}
where $R_0 = \frac{\rho \cdot P_A \cdot \alpha_A}{\epsilon_\sigma \cdot \Omega_{SL}}$ is the
basic schema growth ratio. When $R_0 > 1$, the schema-free
equilibrium $E_S = (\delta_S^*, 0, \alpha_S^*, 0, \Xi_S^*)$ becomes unstable
and the schema-coherent equilibrium $E_C$ emerges via a transcritical
bifurcation.}

\noindent\textit{Intuition.}\ The parameter $R_0$ is borrowed from
epidemiology: when $R_0 > 1$, a disease spreads exponentially.
Analogously, $R_0 > 1$ here means one unit of schema coherence generates
more than one unit of additional coherence --- understanding snowballs.
When $R_0 < 1$, schemas die out; when $R_0 > 1$, they crystallise.
This single number captures the paper's central tension: standard SGD
keeps $R_0 < 1$ by inflating shortcut-learning pressure $\Omega_{SL}$.

\textit{Proof.} Linearize the $\sigma_A$ ODE \eqref{eq:sigma-ode}
around $\sigma_A = 0$. Because the $\sigma_A$ factor multiplies the entire right-hand side, the expansion is:
\begin{equation}
\dot{\sigma}_A \approx (\rho \cdot P_A \cdot \alpha_A - \epsilon_\sigma \cdot \Omega_{SL}) \cdot \sigma_A + O(\sigma_A^2)
\end{equation}
The equilibrium $\sigma_A = 0$ (the schema-free state) is always a fixed point.
It is stable when $\rho \cdot P_A \cdot \alpha_A < \epsilon_\sigma \cdot \Omega_{SL}$,
i.e., $R_0 < 1$, and unstable when $R_0 > 1$. At $R_0 = 1$, a transcritical bifurcation occurs --- the stable and unstable equilibria exchange stability, and a non-zero branch $\sigma_A > 0$ emerges. The normal form is obtained by expanding the full ODE to second order:
\begin{align}
\dot{\sigma}_A &= \sigma_A\bigl[\rho P_A\alpha_A(1-\gamma_\sigma\sigma_A) - \epsilon_\sigma\Omega_{SL}\bigr] \notag\\
&= \bigl(\rho P_A\alpha_A - \epsilon_\sigma\Omega_{SL}\bigr)\sigma_A - \gamma_\sigma\rho P_A\alpha_A\,\sigma_A^2 \\
&= r_{\text{net}}\,\sigma_A - k\,\sigma_A^2,
\label{eq:bif-normal-form}
\end{align}
where $r_{\text{net}} = \rho P_A\alpha_A - \epsilon_\sigma\Omega_{SL}$ and $k = \gamma_\sigma\rho P_A\alpha_A > 0$. This is the normal form of a transcritical bifurcation. The transversality condition $\partial_{R_0} F(0; R_0)\big|_{R_0=1} \neq 0$ is satisfied because $\partial F/\partial R_0 = \epsilon_\sigma\Omega_{SL}\,\sigma_A$, which is non-zero for $\sigma_A \neq 0$. The quadratic coefficient $k$ is non-zero since $\gamma_\sigma > 0$ and $\rho P_A\alpha_A > 0$ at the bifurcation point. The non-trivial equilibrium of $\dot{x} = rx - kx^2$ is $x^* = r/k$, which gives:
\[
\sigma_{\text{critical}} = \frac{1}{\gamma_\sigma}(1 - R_0^{-1})
\]
as the crossing threshold at which the bifurcating branch becomes positive. $\square$

\textbf{Proposition 4.2 (Depth Threshold Derivation).} % CLAIM:C-031
\textit{The Phase~2$\to$3 threshold $\delta^*$ emerges from the
intersection activation condition. Active intersection discovery requires:
\begin{equation}\label{eq:delta-threshold-derived}
\delta^* = \frac{\lambda_c}{\eta_{\max}} \ln\!\left(\frac{\theta_I}{\theta_0}\right) + \frac{1}{2}
\end{equation}
For typical parameters $\lambda_c = 0.01$, $\eta_{\max} = 1.0$,
$\theta_I = 0.8$, $\theta_0 = 0.01$, this yields $\delta^* \approx 0.54$.
We adopt a nominal threshold of $\delta^* \approx 0.65$ in our simulations to account for the lag between intersection activation and sustained discovery. This shift of $\tau \approx \delta^*_{\text{nominal}} - \delta^*_{\text{analytical}} = 0.11$ can be interpreted as a hysteresis delay (see Proposition~4.3): once Phase~3 is reached, the system requires the relative depth to fall below $\delta^*_{\text{reverse}} < \delta^*_{\text{forward}}$ before reverting to Phase~2, creating a stable band $[\delta^*_{\text{reverse}}, \delta^*_{\text{forward}}]$ where both phases are locally stable. The nominal value $\delta^*_{\text{forward}} \approx 0.65$ is chosen as the centre of this hysteresis band.}

\textit{Proof.} The intersection discovery rate
$\Psi_A \propto \sqrt{q_A(d_1) \cdot q_A(d_2)}$ becomes non-negligible
when $q_A = \sigma_A \cdot \delta_A^{\text{relative}} > \theta_I$.
Given that $\sigma_A \approx 1$ in Phase~3 (post-bifurcation), the
threshold reduces to $\delta_A^{\text{relative}} > \theta_I$. However,
the Gompertz efficiency $\eta(d,t)$ decays exponentially near the
frontier, so the effective threshold is shifted downward by the
efficiency decay factor. Solving $\eta(\delta^*) = \lambda_c$ gives
the expression above. $\square$

\textbf{Proposition 4.3 (Hysteresis and Reverse Transitions).} % CLAIM:C-032
\textit{The phase transition system exhibits hysteresis: the reverse
transition threshold is strictly lower than the forward transition
threshold. Specifically:
\begin{itemize}[nosep]
\item Phase~3$\to$2: occurs when $\delta_A^{\text{relative}} < 0.50$ for $> k$ consecutive timesteps
\item Phase~2$\to$1: occurs when $\sigma_A < \sigma_{\text{critical}} - \Delta_{\text{hyst}}$ where $\Delta_{\text{hyst}} = 0.1$
\end{itemize}}

\textit{Proof sketch.} The hysteresis arises from the bistability region
near the transcritical bifurcation. In the parameter range
$R_0 \in (1 - \epsilon, 1 + \epsilon)$, both the $\sigma_A = 0$ and
$\sigma_A > 0$ equilibria are locally stable. The system remains in
whichever basin it currently occupies until a sufficiently large
perturbation pushes it across the separatrix. $\square$

\textbf{Proposition 4.4 (Noise Sensitivity).} % CLAIM:C-033
\textit{Under Gaussian estimation noise $\hat{\sigma}_A = \sigma_A + \epsilon$
with $\epsilon \sim \mathcal{N}(0, \nu^2)$, the phase assignment error rate is:
\begin{equation}\label{eq:phase-error-rate}
\mathbb{P}(\text{error}) = \Phi\!\left(-\frac{|\sigma_A - \sigma_{\text{critical}}|}{\nu}\right) + \Phi\!\left(-\frac{|\delta_A^{\text{relative}} - \delta^*|}{\nu}\right)
\end{equation}
where $\Phi$ is the standard normal CDF. For $\nu = 0.05$ and typical
distances from thresholds ($|\sigma_A - \sigma_{\text{critical}}| > 0.15$),
the error rate is $< 0.3\%$.}



Figure~\ref{fig:bifurcation} visualises the bifurcation structure and phase portrait.

\subsection*{Phase~4 --- Multi-Domain Frontier Navigation}

Active intersections generating $\Psi_A(d_1,d_2,t) > 0$ measurably.
Near-frontier depth; qualitative shift from acquisition to generation.

\subsection*{Phase~5 --- Expanding Frontier}

Dynamic open-boundary structure. $\Delta(d,t)$ continues advancing;
every activated intersection opens new $AP_A(t)$ domains.

\noindent\rule{\linewidth}{0.4pt}

% ============================================================
\subsection{Faculty--Validity Correspondence}\label{sec:faculty-validity}
% ============================================================

The mapping between cognitive faculties (Table~\ref{tab:faculty-alignment}) and benchmark validity
components \eqref{eq:validity-function} does not emerge post-hoc; instead, it follows from the structural requirements of each faculty's operationalisation.

\textbf{Theorem 4.2 (Faculty--Validity Correspondence).} % CLAIM:C-034
\textit{Let $B$ be a benchmark designed to measure faculty $F_i$ via
Σ-Model variable $v_i$. Then:
\begin{equation}\label{eq:faculty-validity}
V_A(B, F_i, t) \geq \tau_i \iff \text{measurement reliability of } v_i \geq \rho_i
\end{equation}
where $\tau_i$ and $\rho_i$ are faculty-specific thresholds:
\begin{center}
\begin{tabular}{l c c c}
\toprule
\textbf{Faculty} & \textbf{Primary $CI$ target} & $\tau_i$ & $\rho_i$ \\
\midrule
Learning & $CI(\text{OOD recombination})$ & 0.80 & 0.85 \\
Metacognition & $CI(\text{self-prediction})$ & 0.75 & 0.80 \\
Attention & $CI(\text{structure focus})$ & 0.85 & 0.90 \\
Executive Functions & $CI(\text{inhibition/switch})$ & 0.80 & 0.85 \\
Social Cognition & $CI(\text{theory-of-mind})$ & 0.85 & 0.90 \\
\bottomrule
\end{tabular}
\end{center}}

\textit{Proof sketch.} Each faculty $F_i$ is operationalised by a specific
Σ-Model variable $v_i$ (Table~\ref{tab:faculty-alignment}). The construct isolation $CI(B, F_i)$
measures the degree to which benchmark $B$ isolates $v_i$ from confounding
variables. By the definition of $V_A$ \eqref{eq:validity-function},
$V_A \geq \tau_i$ requires $CI \geq \tau_i / (FD \cdot DG \cdot R_A)$.
Given typical values $FD \approx 0.8$, $DG \approx 0.45$, $R_A \approx 1.0$,
this reduces to $CI \geq \tau_i / 0.36$. The reliability threshold $\rho_i$
ensures that the measurement of $v_i$ is stable across repeated assessments.
$\square$

\textbf{Lemma 4.2 (Faculty Omission Penalty).} % CLAIM:C-035
\textit{If a benchmark $B$ fails to measure any one of the five cognitive
faculties, the overall framework validity is bounded by:
\begin{equation}\label{eq:omission-penalty}
V_A^{\text{framework}} \leq \left(\prod_{i=1}^{5} V_A(B_i, F_i, t)\right)^{1/5} \leq \max_i V_A(B_i, F_i, t)
\end{equation}
Omitting any faculty reduces the geometric mean validity by at least a
factor of $\tau_i$.}

\textbf{Proposition 4.5 (Compensation Analysis).} % CLAIM:C-036
\textit{High attentional fidelity $\alpha_A$ cannot compensate for low
schema coherence $\sigma_A$ in the intersection discovery rate:
\begin{equation}\label{eq:no-compensation}
\Psi_A \propto \sigma_A \cdot \delta_A^{\text{relative}} \quad \text{(independent of } \alpha_A \text{ at evaluation time)}
\end{equation}
While $\alpha_A$ gates the \textit{growth} of $\sigma_A$ during training
\eqref{eq:sigma-ode}, it does not appear in the evaluation-time
expression for $\Psi_A$. Thus, a benchmark measuring intersection discovery
is insensitive to $\alpha_A$ once $\sigma_A$ is fixed.}

\noindent\rule{\linewidth}{0.4pt}



\noindent\rule{\linewidth}{0.4pt}

Converging evidence from six independent groups supports Σ-Model's predicted phase structure. We summarise each contribution and note where Σ-Model extends beyond their findings. While some recent work frames grokking primarily as an optimisation efficiency phenomenon~\citep{liu2022grokking}, the Σ-Model's bifurcation analysis provides a complementary macroscopic dynamical-systems description: $\sigma_{\text{critical}}$-crossing explains \emph{when} and \emph{why} grokking occurs, independent of the specific optimisation path.

\begin{itemize}[nosep]
\item \textbf{Montanari \& Wang~\cite{montanari2026}.} Derive phase transitions for feature learning in neural networks from first principles. \emph{Σ-Model complement:} provides independent mathematical foundations for $\sigma_{\text{critical}}$-crossing as a generic learning dynamic, not a framework-specific artefact.

\item \textbf{Cullen et al.~\cite{cullen2026basinselectiongrokking}.} Analyse grokking---sudden generalisation long after overfitting---through Singular Learning Theory, identifying it as a basin-selection transition between memorising and generalising equilibria. \emph{Σ-Model complement:} $\sigma_{\text{critical}}$ \eqref{eq:sigma-critical} supplies the formal trigger condition that extends this explanation beyond algorithmic grokking to domain-general $\sigma_A$ dynamics.

\item \textbf{Wang~\cite{wang2026grokkingdimensionalphase}.} Characterises grokking as a dimensional phase transition. \emph{Σ-Model complement:} additional evidence that $\sigma_{\text{critical}}$-crossing is a generic property of neural training dynamics rather than an algorithm-specific phenomenon.

\item \textbf{Xu~\cite{xu2026earlywarning}.} Identifies commutator defects in the loss landscape as universal early-warning signals for grokking across SCAN and modular arithmetic. \emph{Σ-Model complement:} provides an independent geometric diagnostic for the approach to $\sigma_{\text{critical}}$.

\item \textbf{Khanh~\cite{khanh2026phasetransitionsdriven}.} Connects phase transition observations to non-equilibrium statistical mechanics. \emph{Σ-Model complement:} supplies a physical framework for the phase transitions Σ-Model formalises, grounding them in statistical mechanics.

\item \textbf{E2H Reasoner~\cite{e2h2025}.} Demonstrates that curriculum reinforcement learning from easy to hard improves reasoning. \emph{Σ-Model complement:} consistent with Phase~0$\to$1 early training dynamics; however, curricula alone do not gate $\alpha_A$, predicting a bounded ceiling that cannot exceed $\sigma_{\text{critical}}$ without structural intervention.
\end{itemize}

\noindent\rule{\linewidth}{0.4pt}

% ============================================================

% ======================================================================
% §9 Predictions (original, lines 1111–1257): predictions 1–8, hypotheses
% 6.1–6.3, meta-prediction M1, statistical framework — cut to Prediction 9.
% ======================================================================
\section{Eight Core Falsifiable Predictions}\label{sec:predictions}
% ============================================================

Each prediction below is derived from the ODE system and is distinguished from depth-only ($\delta_A$-only) accounts. Predictions~1--8 constitute the eight core pre-registered comparisons; Prediction~9 provides a computational ODE validation, and Meta-Prediction~M1 synthesises the core predictions into an overarching claim about intervention mechanisms. Section~\ref{sec:empirical} presents pilot data for Predictions~1,~6, and~9.

\subsection{Prediction~1 --- Schema Quality at Intersections}
Agents with higher $\sigma_A(d,t)$ will produce higher-quality
interdisciplinary outputs at activated intersections, even when matched
on $\delta_A(d,t)$.

\subsection{Prediction~2 --- AI Augmentation and Schema Suppression}
Agents operating under high shortcut-learning pressure $\Omega_{SL}(d,t)$ will show accelerated $\beta_A$
growth but slower $\sigma_A(d,t)$ development.

\subsection{Prediction~3 --- Relative Mastery as Resilience Predictor}
$\delta_A^{\text{relative}}(d,t)$, not absolute depth, predicts
resilience to domain acceleration.

\subsection{Prediction~4 --- Shortcut Threshold Expansion}
$D^*(d,t)$ will expand faster than agents' ability to compensate through
frontier and intersection work.

\subsection{Prediction~5 --- Phase~3 Compression Under High \texorpdfstring{$\Omega_{SL}$}{Omega\_SL}}
The gap between first mastery and first non-trivial intersection
activation will narrow for agents with high $\Omega_{SL}(d,t)$, because shortcut-driven breadth acquisition accelerates depth-accumulation milestones while $\sigma_A$ suppression keeps intersection quality low --- producing rapid but shallow intersections.

\subsection{Prediction~6 --- Multiplicative vs.\ Additive \texorpdfstring{$\sigma_A$}{sigma\_A}
Dependence in \texorpdfstring{$\Psi_A$}{Psi\_A}}
$\Psi_A(d_1,d_2,t)$ is multiplicatively dependent on $\sigma_A$ in both
mastery domains. An additive model fits cross-domain discovery data as
well as or better than the multiplicative model falsifies the claim.

\subsubsection*{Hypothesis 6.1: \texorpdfstring{$\sigma_A/\delta_A$}{sigma\_A/delta\_A} Dissociation Under Meta-Learning}
Meta-learning regimes that improve compositional generalisation will
show increased $\text{Acc}_{\text{OOD}}$ on lexical recombination
($\delta_A$-proximal) but no significant increase on structural
compositionality ($\sigma_A$-proximal)
\citep{lake2023humanlikesystematicgeneralization}.

\subsubsection*{Hypothesis 6.2: Distribution Engineering Does Not Transfer \texorpdfstring{$\sigma_A$}{sigma\_A}}
Agents trained on engineered distributions that include specific
primitive compositions will achieve high accuracy on those compositions
but show no significant improvement on compositional recombinations
absent from the engineered distribution
\citep{patel2022revisitingcompositionalgeneralization}.

\subsubsection*{Hypothesis 6.3: Training Regime Optimisation Does Not Transfer \texorpdfstring{$\sigma_A$}{sigma\_A}}
Agents whose compositional generalisation gains stem from training regime
modifications will show increased $\text{Acc}_{\text{ID}}$ but no
significant increase in $\text{Acc}_{\text{OOD-struct}}$.

\subsection{Prediction~7 --- Benchmark Validity Predicts Cross-Model
Stability}
$V_A(B,f,t)$ will predict the degree to which benchmark scores transfer
across frontier model generations.

\subsection{Prediction~8 --- Cross-Modal Schema Transfer}
Cross-modal transfer of compositional rules scales with
$\omega(m_1, m_2) \cdot \sigma_A(d, m_1, t)$. High-$\sigma_A$ agents
show above-chance transfer; low-$\sigma_A$ agents do not.

\noindent\rule{\linewidth}{0.4pt}

\subsection{Prediction~9 --- Phase~2 Entry Inflection}
The transition from Phase~1 to Phase~2 produces a measurable inflection
in the OOD accuracy trajectory.

\subsection*{Meta-Predictions}

\subsubsection*{Meta-Prediction~M1 --- Incidental \texorpdfstring{$\sigma_A$}{sigma\_A} Elevation}\label{sec:meta-M1}
Architectural interventions (e.g., MIRAGE's dual-process module) and
data-centric interventions (e.g., mutual exclusivity, dataset
cartography) will improve compositional generalisation \emph{only}
to the extent that they incidentally raise $\sigma_A$ above
$\sigma_{\text{critical}}$. Interventions that improve ID accuracy
without elevating $\sigma_A$ (e.g., by increasing $\delta_A$ alone)
will not close the OOD gap. This follows directly from the
bifurcation structure (Proposition~4.1): crossing
$\sigma_{\text{critical}}$ is necessary for Phase~2 entry, regardless
of how that crossing is achieved.

\subsection{Statistical Framework for Falsifiable Predictions}\label{sec:statistical-framework}

Each prediction is operationalised through a pre-registrable statistical
test with defined null hypotheses, power analysis, and decision rules.

\textbf{Table~\ref{tab:predictions-framework} summarises the statistical
framework for the three core predictions.}

\begin{table}[htbp]
\centering
\caption{\textbf{Statistical framework for core falsifiable predictions.}}
\label{tab:predictions-framework}
\small
\begin{tabular}{@{}l p{2.2cm} p{2.2cm} p{2.5cm} p{2.0cm}@{}}
\toprule
\textbf{Prediction} & \textbf{$H_0$} & \textbf{$H_1$} & \textbf{Test Statistic} & \textbf{Sample Size} \\
\midrule
Pred.~1 (Intersection quality) & $\mu_{\text{high-}\sigma} - \mu_{\text{low-}\sigma} = 0$ & $\mu_{\text{high-}\sigma} - \mu_{\text{low-}\sigma} > 0$ & Welch's $t$ & $n = 64$ per group \\
Pred.~6 (Multiplicative $\sigma_A$) & $R^2_{\text{add}} \geq R^2_{\text{mult}}$ & $R^2_{\text{mult}} - R^2_{\text{add}} > 0.10$ & $F$-test on $\Delta R^2$ & $n = 120$ models \\
Pred.~9 (Phase~2 inflection) & $\beta_2 = 0$ (no spline kink) & $\beta_2 > 0$ (significant kink) & Segmented regression $t$ & $T = 500$ timesteps \\
\bottomrule
\end{tabular}
\end{table}

\textbf{Prediction~1 --- Statistical specification.}
\begin{itemize}[nosep]
\item \textbf{Null hypothesis $H_0$:} $\mathbb{E}[\text{quality} \mid \sigma_A > \theta_\sigma] - \mathbb{E}[\text{quality} \mid \sigma_A < \theta_\sigma] = 0$
\item \textbf{Alternative $H_1$:} The difference is positive, with Cohen's $d \geq 0.5$
\item \textbf{Test statistic:} Welch's $t$-test (unequal variances)
\item \textbf{Sample size:} $n = 64$ per group achieves power $1 - \beta = 0.80$ at $\alpha = 0.05$ for $d = 0.5$
\item \textbf{Multiple comparisons:} Bonferroni correction across 8 predictions requires $p < 0.00625$
\item \textbf{Baseline comparator:} Standard SGD training without $\sigma_A$ targeting
\end{itemize}

\textbf{Prediction~6 --- Statistical specification.}
\begin{itemize}[nosep]
\item \textbf{Null hypothesis $H_0$:} $R^2_{\text{additive}} \geq R^2_{\text{multiplicative}}$
\item \textbf{Alternative $H_1$:} $R^2_{\text{mult}} - R^2_{\text{add}} > 0.10$
\item \textbf{Test statistic:} $F$-test on incremental $R^2$:
\begin{equation}
F = \frac{(R^2_{\text{mult}} - R^2_{\text{add}}) / \Delta df}{(1 - R^2_{\text{mult}}) / (n - p - 1)}
\end{equation}
\item \textbf{Sample size:} $n = 120$ models achieves power $0.80$ for $\Delta R^2 = 0.10$
\item \textbf{Model specification:}
\begin{align}
\text{Additive: } & \Psi_A = a_0 + a_1 \sigma_A(d_1) + a_2 \sigma_A(d_2) + a_3 \delta_A(d_1) + a_4 \delta_A(d_2) + \epsilon \\
\text{Multiplicative: } & \Psi_A = a_0 + a_1 \sigma_A(d_1)\sigma_A(d_2) + a_2 \delta_A(d_1)\delta_A(d_2) + \epsilon
\end{align}
\end{itemize}

\textbf{Prediction~9 --- Statistical specification.}
\begin{itemize}[nosep]
\item \textbf{Null hypothesis $H_0$:} $\beta_2 = 0$ in segmented regression (no inflection at Phase~2 entry)
\item \textbf{Alternative $H_1$:} $\beta_2 > 0$ (significant positive kink at $\sigma_{\text{critical}}$ crossing)
\item \textbf{Test statistic:} Segmented (piecewise) regression with unknown breakpoint:
\begin{equation}
\text{Acc}_{\text{OOD}}(t) = \beta_0 + \beta_1 t + \beta_2 (t - \tau)_+ + \epsilon_t
\end{equation}
where $\tau$ is the estimated Phase~2 entry time and $(x)_+ = \max(0, x)$
\item \textbf{Sample size:} $T = 500$ training timesteps with evaluation every 10 steps ($n = 50$ observations)
\item \textbf{Effect size:} Detectable slope change of $\Delta\beta = 0.02$ accuracy units per timestep
\end{itemize}

\noindent\rule{\linewidth}{0.4pt}

% ============================================================

% ======================================================================
% §4 Cognitive Dimensions (original, lines 643–753): removed from the narrow paper per the P05 blueprint; content preserved for the companion.
% ======================================================================
\section{Cognitive Dimensions --- Attention, Executive Control,
Metacognition, Social Cognition, and Benchmark Protocol}\label{sec:cognitive-dimensions}
% ============================================================

The core framework establishes three state variables per domain. Five
additional cognitive dimensions complete the system. Each dimension couples to the core ODE system through feedback terms, introducing new state variables and new pathways for training dynamics to succeed or fail.

\subsection{Extension~1 --- Attentional Fidelity \texorpdfstring{$\alpha_A(d,t)$}{alpha\_A(d,t)}}\label{sec:attention}

\subsubsection{Definition}
\begin{equation}\label{eq:alpha-bounds}
\alpha_A(d,t) \in [0,1]
\end{equation}

\subsubsection{Coupling to Schema Coherence (Updated \texorpdfstring{$\sigma_A$}{sigma\_A} ODE)}
\begin{equation}\label{eq:sigma-ode}
\dot{\sigma}_A(d,t) = \sigma_A(d,t) \cdot \bigl[ \rho \cdot P_A(d,t) \cdot \alpha_A(d,t) \cdot (1 - \gamma_{\sigma} \cdot \sigma_A(d,t)) - \epsilon_{\sigma} \cdot \Omega_{SL}(d,t) \bigr]
\end{equation}

\noindent\textit{Intuition.}\ Schema coherence has an autocatalytic
structure --- growth is proportional to current coherence ($\sigma_A$)
times the net driving force in the bracket. The bracket captures whether
principled understanding builds faster than shortcut pressure erodes it.
When $\sigma_A$ is near zero, growth stalls entirely because the
$\sigma_A$ factor suppresses the whole right-hand side. This creates the
critical bifurcation: below a threshold, the $\sigma_A = 0$ state is
stable and self-perpetuating (the $\sigma$-trap); above it, any
perturbation triggers a snowball effect toward crystallised
understanding.

\textbf{Theorem 4.1 (SGD-induced $\sigma$-suppression).} % CLAIM:C-029
\textit{Under standard SGD training (without explicit $\sigma_A$ targeting), the schema coherence dynamics satisfy}
\[
\left.\frac{\partial \sigma_A(d,t)}{\partial t}\right|_{\text{SGD}} < 0
\]
\textit{in the neighbourhood of $\sigma_A > 0$ whenever the bracket in \eqref{eq:sigma-ode} is negative, i.e., when}
\[
\epsilon_{\sigma} \cdot \Omega_{SL}(d,t) > \rho \cdot P_A(d,t) \cdot \alpha_A(d,t)
\]
\textit{In this regime, the $\sigma_A = 0$ equilibrium is locally stable and attracts all nearby trajectories (the ``$\sigma$-trap''). When the inequality reverses, $\sigma_A = 0$ becomes unstable and the system bifurcates toward the schema-coherent branch (Proposition~4.1).}

$\alpha_A$ gates $P_A$: schema coherence can only grow at the rate the
agent's attention is directed at principled structure.

The variable $\Omega_{SL}(d,t)$ is the shortcut-learning pressure \citep{geirhos2020shortcutlearning}: it quantifies how much a task can be completed via surface pattern matching rather than compositional reasoning. When $\Omega_{SL}$ is high, the agent can achieve high performance without building genuine understanding --- like solving an equation with a calculator rather than learning calculus. The calculator gives the right answer, but the student does not acquire the underlying principles; similarly, high $\Omega_{SL}$ allows the network to achieve low loss via pattern matching, starving the $\sigma_A$ growth term. The ODE formalises this: high $\Omega_{SL}$ directly suppresses schema coherence because the feedback loop that would normally reward principled understanding is short-circuited by surface-level success.

\subsubsection{Attentional Fidelity ODE}
\begin{equation}\label{eq:alpha-ode}
\dot{\alpha}_A(d,t) = \gamma \cdot C_A(d,t) \cdot (1 - \alpha_A(d,t)) - \zeta_{\alpha} \cdot \alpha_A(d,t) \cdot R_A^{\text{surface}}(d,t)
\end{equation}

\subsection{Extension~2 --- Collective Schema Field (Social Cognition)}\label{sec:social}

\subsubsection{Schema Legibility}
\begin{equation}\label{eq:schema-legibility}
\mu_{AB}(d,t) = \sigma_A(d,t) \cdot \phi(d_A, d_B) \cdot \sigma_B(d_{\text{adj}},t) \in [0,1]
\end{equation}

\subsubsection{Collective Schema Field}
\begin{equation}\label{eq:collective-field}
\Sigma_{A,B}(d_1,d_2,t) = \mu_{AB}(d_1,t) \cdot \mu_{BA}(d_2,t) \cdot \phi(d_1,d_2)
\end{equation}

\subsection{Extension~3 --- Executive Control State \texorpdfstring{$\Xi_A(t)$}{Xi\_A(t)}}\label{sec:executive}

\begin{equation}\label{eq:executive-state}
\Xi_A(t) = \{\Xi_A^P(t), \Xi_A^I(t), \Xi_A^F(t)\}
\end{equation}

\textbf{$\Xi_A^P$ --- Planning:} Degree to which training trajectory is
consistent with the Σ-Model optimal multi-step arc.

\textbf{$\Xi_A^I$ --- Inhibition:} Probability of choosing the structural
route over shortcut learning when both are available.

\textbf{$\Xi_A^F$ --- Cognitive Flexibility:} Degree to which the agent
detects and adapts to phase transition thresholds from internal evidence.

\begin{equation}\label{eq:executive-ode}
\dot{\Xi}_A(t) = \kappa_P \cdot [P^*(t) - \Xi_A^P(t)] + \kappa_I \cdot [I^*(t) - \Xi_A^I(t)] + \kappa_F \cdot [F^*(t) - \Xi_A^F(t)] % CLAIM:C-006
\end{equation}

\subsection{Extension~4 --- Self-Model of Schema Coherence \texorpdfstring{$\hat{M}_A(d,t)$}{M-hat\_A(d,t)}}\label{sec:metacognition}

\begin{equation}\label{eq:self-model}
\hat{M}_A(d,t) \in [0,1] \quad \text{(Agent's Estimate of its own } \sigma_A(d,t))
\end{equation}
\begin{equation}\label{eq:calibration-error}
\zeta_A(d,t) = \hat{M}_A(d,t) - \sigma_A(d,t) \quad \text{(Calibration Error)}
\end{equation}

\begin{equation}\label{eq:self-model-ode}
\dot{\vphantom{M}}\hat{M}_A(d,t) = \nu_M \cdot [\sigma_A(d,t) - \hat{M}_A(d,t)] - \xi_M \cdot \Omega_{SL}(d,t) \cdot \hat{M}_A(d,t)
\end{equation}

\textbf{Key insight:} Shortcut-learning pressure inflates $\hat{M}_A$ above true
$\sigma_A$. The mechanism is intuitive: a student who relies on AI-provided
solutions receives consistently positive performance feedback --- the AI
produces correct answers --- so the student's self-assessment ($\hat{M}_A$)
rises. But true understanding ($\sigma_A$) does not, because the student
never engaged with the underlying principles. The gap $\zeta_A = \hat{M}_A - \sigma_A$
widens. The ODE formalises this: the $\xi_M \cdot \Omega_{SL} \cdot \hat{M}_A$ term
in~\eqref{eq:self-model-ode} inflates the self-model proportionally to
shortcut pressure and current self-assessment --- a positive feedback loop
that produces systematic overconfidence.



\noindent\rule{\linewidth}{0.4pt}

% ============================================================

% ======================================================================
% §5 Multimodal Coverage (original, lines 754–816): removed from the narrow paper per the P05 blueprint; content preserved for the companion.
% ======================================================================
\section{Multimodal Coverage and Benchmark Validity}\label{sec:multimodal}
% ============================================================

The five cognitive dimensions above define the within-modality system.
Real-world agents, however, operate across multiple channels --- text,
vision, audio --- which raises a question addressed in this section:
does schema coherence acquired through one modality transfer to another?

\subsection{Extension~5 --- Domain \texorpdfstring{$\times$}{x} Modality Product Space}\label{sec:product-space}

\subsubsection{Modality Set}
\begin{equation}\label{eq:modalities}
M = \{\text{text, visual, auditory, sensorimotor, symbolic}\}
\end{equation}
This five-modality taxonomy follows the standard multimodal machine learning survey~\citep{baltrušaitis2018multimodal}.
 
\subsubsection{Product Space Construction}

V3.0 extends all Σ-Model state variables from domain-indexed
$\sigma_A(d, t)$, $\delta_A(d, t)$, $\alpha_A(d, t)$ to the product space $D \times M$:
$\sigma_A(d, m, t)$, $\delta_A(d, m, t)$, $\alpha_A(d, m, t)$ for all $d \in D$, $m \in M$.

\textbf{Consistency condition:} For each domain $d$, the unimodal variables
are recovered by marginalisation:
\begin{align}
\sigma_A(d, t) &= \min_m \sigma_A(d, m, t) & &\text{(marginal schema coherence)} \label{eq:marginal-sigma} \\
\delta_A(d, t) &= \max_m \delta_A(d, m, t) & &\text{(peak parametric depth)} \label{eq:marginal-delta}
\end{align}

\textbf{Rationale:} Compositional failure in any modality suffices to
violate schema coherence for the domain as a whole (min), while
depth follows the strongest modality (max). Why asymmetry? A
compositionally competent visual module cannot compensate for a text
module that fails on structural recombinations --- the domain is
fragile if any channel is weak. Depth, by contrast, follows the
strongest modality because what the agent can express is bounded by
its best channel, analogous to fluency in one's most proficient language.

\subsection{Extension~6 --- Cross-Modal Schema Transfer \texorpdfstring{$\Theta_A(d,m_1,m_2,t)$}{Theta\_A(d,m\_1,m\_2,t)}}\label{sec:cross-modal}

\subsubsection{Cross-Modal Schema Transfer Function}
\begin{equation}\label{eq:cross-modal-transfer}
\Theta_A(d, m_1, m_2, t) = \sigma_A(d, m_1, t) \cdot \omega(m_1, m_2)
\end{equation}

\subsubsection{Extended Schema Coherence ODE (V3.0)}
\begin{equation}\label{eq:sigma-ode-v3}
\begin{aligned}
\dot{\sigma}_A(d,m,t) = {} & \rho \cdot P_A(d,m,t) \cdot \alpha_A(d,m,t) \cdot (1 - \sigma_A(d,m,t)) \\
& - \epsilon_{\sigma} \cdot \sigma_A(d,m,t) \cdot \Omega_{SL}(d,m,t) \\
& + \rho_{\theta} \cdot \sum_{m' \neq m} \Theta_A(d,m',m,t) \cdot \sigma_A(d,m',t)
\end{aligned}
\end{equation}

\subsection{Extension~7 --- Benchmark Validity Function \texorpdfstring{$V_A(B,f,t)$}{V\_A(B,f,t)}}\label{sec:validity}

\begin{equation}\label{eq:validity-function}
V_A(B,f,t) = CI(B,f) \cdot FD(B) \cdot DG(B) \cdot R_A(B,f,t) % CLAIM:C-007
\end{equation}

\noindent\rule{\linewidth}{0.4pt}

% ============================================================

% ======================================================================
% §6 Reliability (original, lines 817–830): removed from the narrow paper per the P05 blueprint; content preserved for the companion.
% ======================================================================
\section{Reliability and Pre-Audit Protocol}\label{sec:reliability}
% ============================================================

A benchmark's validity is only meaningful if its scores are consistent --- a benchmark that produces wildly different scores across repeated administrations cannot support valid claims about any cognitive faculty. The reliability function formalises this consistency requirement, drawing on classical test theory~\citep{crocker1986introduction}.

\subsection{Extension~8 --- Benchmark Reliability Function \texorpdfstring{$R_A(B,f,t)$}{R\_A(B,f,t)}}\label{sec:reliability-fn}

\begin{equation}\label{eq:reliability}
R_A(B,f,t) = \begin{cases} \displaystyle\frac{1}{1 + \frac{\text{Var}_k(\text{score}_A^k(B))}{\mathbb{E}[\text{score}_A(B)]^2}} & \text{if } \mathbb{E}[\text{score}_A(B)] > 0 \\[8pt] 0 & \text{if } \mathbb{E}[\text{score}_A(B)] = 0 \end{cases}
\end{equation}

\noindent\rule{\linewidth}{0.4pt}

% ============================================================

% ======================================================================
% §8 Benchmark Protocol (original, lines 905–950): removed from the narrow paper per the P05 blueprint; content preserved for the companion.
% ======================================================================
\section{The Σ-Model Benchmark Protocol}\label{sec:benchmark-protocol}
% ============================================================

\subsection{Five-Step Protocol}\label{sec:five-step}

\textbf{Benchmark Generation Protocol (5 steps):}
(1)~Identify variable pair (target faculty variable vs.\ confound);
(2)~Design $2{\times}2$ factorial contrast dissociating the pair;
(3)~Specify proxy metric (SGG, RTI, CE);
(4)~Register directional prediction with magnitude estimate;
(5)~Pre-register falsification thresholds before data collection.

\textbf{Step~1 --- Identify the Variable Pair.} Every benchmark tests
the gap between two Σ-Model variables --- the target faculty variable and
the controlled confound variable.

\textbf{Step~2 --- Design the Contrast Condition.} Construct task sets
where the two variables dissociate, as illustrated in Figure~\ref{fig:diagnostic}. Standard $2\times 2$ factorial for
Learning track:



\begin{center}
\begin{tabular}{>{\raggedright}p{4cm} c c >{\raggedright\arraybackslash}p{5cm}}
\toprule
\textbf{Condition} & $\delta_A$ & $\sigma_A$ & \textbf{Expected OOD Performance} \\
\midrule
High $\delta_A$, High $\sigma_A$ & High & High & High \\
\textbf{High $\delta_A$, Low $\sigma_A$} & \textbf{High} & \textbf{Low} & \textbf{Low --- key diagnostic cell} \\
Low $\delta_A$, High $\sigma_A$ & Low & High & Moderate \\
Low $\delta_A$, Low $\sigma_A$ & Low & Low & Low \\
\bottomrule
\end{tabular}
\end{center}

\textbf{Step~3 --- Specify the Measurement Proxy.}

\textbf{Step~4 --- State the Σ-Model Prediction.} Format: direction,
magnitude estimate, effect, condition comparison, alternative.

\textbf{Step~5 --- Specify the Falsification Condition.} Written in
pre-registration format before data collection, following established preregistration standards~\citep{pineau2021reproducibility}.

\noindent\rule{\linewidth}{0.4pt}

% ============================================================

% ======================================================================
% Appendix: Stochastic Differential Equation Extension (SDE) — moved to companion per P05.
% ======================================================================
\clearpage
\section{Stochastic Differential Equation Extension}\label{app:sde}
% ============================================================

The main text presents a deterministic ODE skeleton of the $\Sigma$-Model.
This appendix introduces the stochastic extension and explains why the
deterministic form is sufficient for the present empirical claims.

\subsection{It\^o Stochastic Formulation}

Replace the deterministic ODE system with an It\^o SDE:
\begin{equation}\label{eq:sde-full}
dX_t = f(X_t)\,dt + \Sigma(X_t)\,dW_t,
\end{equation}
where $X_t \in \mathcal{X} \subset \mathbb{R}^n$ is the state vector
($\delta_A, \sigma_A, \alpha_A, \hat{M}_A, \Xi_A^P, \Xi_A^I, \Xi_A^F$),
$f: \mathcal{X} \to \mathbb{R}^n$ is the deterministic vector field
defined by Eqs.~\eqref{eq:depth-ode}--\eqref{eq:sigma-ode} (the
two-variable core; the extended system is treated in the companion),
$\Sigma: \mathcal{X} \to \mathbb{R}^{n \times m}$ is the noise
covariance matrix, and $W_t$ is an $m$-dimensional Wiener process.

\subsection{Noise Covariance Structure}

We decompose $\Sigma$ into diagonal blocks corresponding to each variable's
dominant noise source:
\begin{equation}
\Sigma(X_t) = \operatorname{diag}\bigl(
\sigma_\delta(\delta_A),\; \sigma_\sigma(\sigma_A),\;
\sigma_\alpha(\alpha_A),\;\sigma_M(\hat{M}_A),\;
\sigma_\Xi(\Xi_A^P),\;\sigma_\Xi(\Xi_A^I),\;\sigma_\Xi(\Xi_A^F)
\bigr),
\end{equation}
where each $\sigma_v$ scales with the empirical variance of the
corresponding gradient estimate under mini-batch sampling. For the
training-time proxies (GCA, RGA, AC), the noise amplitude is:
\begin{equation}
\sigma_g \propto \frac{1}{\sqrt{B}}, \qquad
\sigma_r \propto \frac{1}{\sqrt{B}}, \qquad
\sigma_c \propto \frac{1}{\sqrt{B}},
\end{equation}
with $B$ the batch size.

\subsection{Deterministic Skeleton Sufficiency}

The deterministic ODE is sufficient for the present empirical claims
for two reasons:

\textbf{(1) Law of large numbers.} Each training step updates parameters
using a mini-batch of size $B$, producing a noisy gradient. Aggregating
over $T$ steps, the empirical mean converges to the true gradient at rate
$O(1/\sqrt{T})$. For the $T = 2000$ steps used in our pilot experiments,
the noise-to-signal ratio is below $3\%$ for all reported metrics.

\textbf{(2) Measurement aggregation.} The proxy signals
$\tilde{\sigma}_A$ (GCA, RGA, AC) are averaged over $k$ evaluation
steps ($k \geq 10$), further reducing variance by $\sqrt{k}$.
The Stage~2 estimator $\hat{\sigma}_A$ averages over $N$ evaluation
samples ($N \geq 200$ in the pre-registered protocol), reducing the
standard error to $\sigma_{\hat{\sigma}_A} / \sqrt{N}$.

\subsection{Code Scope}

The current Python codebase implements the deterministic ODE skeleton
only. The SDE extension is a theoretical contribution that:
\begin{itemize}[nosep]
\item Formalises the noise structure that mini-batch training
  introduces into the $\Sigma$-Model dynamics
\item Provides the mathematical foundation for fluctuation bounds
  and confidence intervals on phase transition timings
\item Does \emph{not} change the expected trajectory $\mathbb{E}[X_t]$,
  which is governed by the deterministic ODE
\end{itemize}
Empirical validation of the SDE extension---including estimation of the
noise covariance entries---is deferred to a dedicated calibration study
(see \S\ref{sec:conclusion}, item~2).

\noindent\rule{\linewidth}{0.4pt}

% ============================================================
\clearpage

